\documentclass[11pt,oneside]{report}
% ======================================================================
%  THE TRIADIC MEASUREMENT SUBSTRATE --- COMPLETE CORPUS (Volumes 1 & 2)
%  Aaron Switzer, 2026
%  Single-file build. Compile: pdflatex x3 (third pass settles the TOC).
% ======================================================================
% ======================================================================
%  TMS Volume 1 -- shared preamble
%  Compiles with pdfLaTeX both locally and on Overleaf (no exotic fonts).
% ======================================================================
\usepackage[T1]{fontenc}
\usepackage[utf8]{inputenc}
\usepackage{amsmath}
\usepackage{amssymb}
\usepackage{mathptmx}            % Times text + math (widely available)
\usepackage[scaled=0.90]{helvet} % sans for headings
\usepackage{courier}

\usepackage[letterpaper,margin=1in]{geometry}
\usepackage{amsthm}
\usepackage{mathtools}
\usepackage{bm}
\usepackage{microtype}
\usepackage{enumitem}
\usepackage{booktabs}
\usepackage{array}
\usepackage{longtable}
\usepackage{caption}
\usepackage{xcolor}
\usepackage{titlesec}
\usepackage{fancyhdr}
\usepackage[hidelinks]{hyperref}
\usepackage{tocbibind}

% ---- spacing -----------------------------------------------------------
\setlength{\parindent}{1.2em}
\setlength{\parskip}{0.35em}
\linespread{1.04}
\setlist{leftmargin=1.6em,topsep=0.25em,itemsep=0.15em,parsep=0pt}

% ---- theorem environments ---------------------------------------------
\newtheorem{theorem}{Theorem}[section]
\newtheorem{lemma}[theorem]{Lemma}
\newtheorem{corollary}[theorem]{Corollary}
\newtheorem{proposition}[theorem]{Proposition}

\theoremstyle{definition}
\newtheorem{definition}[theorem]{Definition}
\newtheorem{axiom}{Axiom}
\newtheorem{claim}[theorem]{Claim}
\newtheorem{principle}[theorem]{Principle}
\newtheorem{conjecture}[theorem]{Conjecture}
\newtheorem{prediction}[theorem]{Prediction}
\newtheorem{property}[theorem]{Property}
\newtheorem{postulate}[theorem]{Postulate}

\theoremstyle{remark}
\newtheorem{remark}[theorem]{Remark}
\newtheorem*{remark*}{Remark}

% ---- section heading look (unnumbered; author supplies the numbers) ----
\titleformat{\section}[hang]{\normalfont\large\bfseries\sffamily}{}{0pt}{}
\titleformat{\subsection}[hang]{\normalfont\normalsize\bfseries\sffamily}{}{0pt}{}
\titleformat{\subsubsection}[hang]{\normalfont\normalsize\bfseries\itshape}{}{0pt}{}
\titlespacing*{\section}{0pt}{1.4em}{0.5em}
\titlespacing*{\subsection}{0pt}{1.0em}{0.35em}
\titlespacing*{\subsubsection}{0pt}{0.8em}{0.3em}

% chapter (= each paper) heading
\titleformat{\chapter}[display]
  {\normalfont\sffamily\bfseries}
  {}{0pt}{\Huge}
\titlespacing*{\chapter}{0pt}{10pt}{24pt}

% ---- page-break quality: avoid stranded headings and orphaned/widowed lines ----
% (applies to both volumes via the shared preamble)
\widowpenalty=10000        % no single line at the top of a page
\clubpenalty=10000         % no single line at the bottom of a page
\displaywidowpenalty=10000 % same, around displayed equations
\brokenpenalty=4000        % discourage page breaks right after a hyphenated line
% titlesec: forbid a page break immediately after a heading (heading stranded at page foot)
\usepackage{needspace}
\usepackage{etoolbox}
\pretocmd{\section}{\Needspace*{4\baselineskip}}{}{}
\pretocmd{\subsection}{\Needspace*{3\baselineskip}}{}{}

% ---- running heads -----------------------------------------------------
\pagestyle{fancy}
\fancyhf{}
\fancyhead[L]{\small\itshape The Triadic Measurement Substrate}
\fancyhead[R]{\small\itshape\nouppercase{\rightmark}}
\fancyfoot[C]{\small\thepage}
\renewcommand{\headrulewidth}{0.4pt}
\fancypagestyle{plain}{\fancyhf{}\fancyfoot[C]{\small\thepage}\renewcommand{\headrulewidth}{0pt}}

% ---- abstract / keywords / references list ----------------------------
\newenvironment{paperabstract}
  {\begin{center}\begin{minipage}{0.88\textwidth}\small
   \noindent\textbf{Abstract.}\itshape\space}
  {\end{minipage}\end{center}\vspace{0.4em}}

\newcommand{\keywords}[1]{\noindent{\small\textbf{Keywords:} #1}\vspace{0.4em}}

% per-paper APA reference list (hanging indent), kept as authored
\newenvironment{reflist}
  {\par\small\setlength{\parindent}{0pt}%
   \setlength{\leftskip}{1.6em}\setlength{\parindent}{-1.6em}%
   \setlength{\parskip}{0.4em}%
   \raggedright\hyphenpenalty=50\exhyphenpenalty=50\relax
   \sloppy}
  {\par}

% convenience
\providecommand{\tightlist}{\setlength{\itemsep}{0pt}\setlength{\parskip}{0pt}}
\providecommand{\TMS}{\textsc{tms}}
\hyphenation{con-fir-ma-tion sub-strate}
\begin{document}

\thispagestyle{empty}
\begin{titlepage}
\centering
\vspace*{2cm}
{\sffamily\bfseries\Large The Triadic Measurement Substrate\par}
\vspace{0.6cm}
{\large $\bar{B} = 3\bar{\alpha}$\par}
\vspace{1cm}
{\sffamily\large Volume 1\par}
\vspace{1.4cm}
{\Large \bfseries Paper 3: Computation Emerges: Surround-Driven Repair, Coherence Bistability, and the First Subsystem\par}
\vfill
{\large Aaron Switzer\par}
\vspace{0.4cm}
{\large 2026\par}
\vspace{1.2cm}
{\small\itshape Excerpted from the complete two-volume corpus, \emph{The Triadic Measurement Substrate}. Full corpus and reproducibility package: \url{https://doi.org/10.5281/zenodo.22035422}\par}
\end{titlepage}
\cleardoublepage
\pagenumbering{arabic}

% ---------------- Volume 1 / paper3 ----------------
\chapter*{Paper 3}
\markboth{Paper 3}{Paper 3}
\addcontentsline{toc}{chapter}{Paper 3: Computation Emerges: Surround-Driven Repair, Coherence Bistability, and the First Subsystem}
\begingroup\centering
{\large\itshape Computation Emerges: Surround-Driven Repair, Coherence Bistability, and the First Subsystem\par}\vspace{0.8em}
{\normalsize Aaron Switzer\par}
{\small May 2026\par}
\endgroup\vspace{1.0em}

\noindent{\small\textbf{Keywords:} TMS, Stabilizer Repair, Hebbian Learning, Coherence Bistability, Coarse-Graining, Subsystems, Universal Computation, Computational Hierarchy}\par\vspace{0.6em}

\section*{Status of Results}
\addcontentsline{toc}{section}{Status of Results}

This paper distinguishes four categories of claim, each labeled where it
appears in the text.

\textbf{Derived.} Results that follow from the five axioms, the
Principle of Generative Economy, and the structures established in
Papers 1 and 2 by explicit construction or proof. Surround-driven repair
as a direct consequence of FCC boundary geometry, requiring no extension
of the [[4,2,2]] repair operator (\S\,2.2); the Imprinting Theorem
stating that a region of diameter \ensuremath{L_{R}} converges to the unique
stabilizer-satisfying continuation of its retarded surround history over
imprint depth \ensuremath{d_{R}} = \ensuremath{\sqrt{2}} \ensuremath{\cdot} \ensuremath{L_{R}/\ell_{p}} flops (\S\,2.3); Hebbian co-firing as a
corollary of the imprint mechanism under structured boundary conditions
(\S\,2.4); the dissolution rate of non-T-closed configurations and the
effective error-rate decay \ensuremath{(1/3)^{d}} (\S\,2.5, closing the Paper 2
deferrals); the recursive applicability of the axioms at the
coarse-grained scale \ensuremath{L_{p},} producing the meta-level triadic rule,
meta-level FCC structure, and meta-level [[4,2,2]] stabilizer
dynamics by the same arguments applied at the same axiomatic content
(\S\,IV); the \ensuremath{(C_{D},} \ensuremath{C_{O},} \ensuremath{C_{C})} trichotomy as the unique minimum
sufficient decomposition of any non-trivial information-transforming
configuration (\S\,5.2); the minimum subsystem extent \ensuremath{L_{\mathrm{sub}}} = \ensuremath{3L_{p}} + \ensuremath{\delta}
(\S\,5.3); the existence of a minimum computing subsystem by explicit
construction (\S\,5.4); and the information-generation rate differential of
one novel P-state per meta-flop per minimum subsystem (\S\,5.5) are derived
in this sense.

\textbf{Structural correspondence.} The coherence bistability threshold
\textbar \ensuremath{\Pi_{\psi}}\textbar\ensuremath{*(L_{R})} = 1 \ensuremath{-} \ensuremath{2\sqrt{c_{\mathrm{TMS}}/(6}} \ensuremath{L_{R}))} of \S\,3.3 is a
derivation of the scaling form, not of the exact coefficient. The
functional dependence \ensuremath{p_{\mathrm{error}}} \ensuremath{\propto} \ensuremath{1/\sqrt{L_{R}}} is forced by balancing per-flop
noise injection against per-flop imprinting capacity under FCC
connectivity. The numerical coefficients depend on combinatorial factors
(notably the C(4,2) = 6 upper bound on per-tetrahedral-unit two-error
configurations) and geometric proportionality factors on the imprinting
side, which are stated as worst-case bounds and reserved for refinement
in Paper 5's simulation work. The threshold's qualitative behavior ---
strictly between 0 and 1, monotonic in \ensuremath{L_{R},} producing a bistable phase
structure --- is forced by the derivation. Its quantitative pinning is
not.

\textbf{Predictions / Conjectures.} Falsifiable claims that follow from
the framework but whose closed forms are not derived within this paper.
The hierarchical lightspeed formula \ensuremath{c^{(k)}} (\S\,6.1) gives the structural
form, with the \ensuremath{L_{p}(k)/L_{p}(k-1)} and \ensuremath{T_{p}(k)/T_{p}(k-1)} ratios identified
as combinatorial quantities requiring enumeration of
stabilizer-satisfying T-closed configurations at each level. The refined
\ensuremath{\beta} = f(d, \ensuremath{K_{\psi,\mathrm{local}},} \ensuremath{K_{\psi,\mathrm{surround}},} \ensuremath{K_{\psi,\mathrm{subsystem}})} prediction (\S\,6.2)
specifies the functional dependences; f's explicit form is deferred. The
spatial-distribution prediction for computational activity in old
substrate regions (\S\,6.3) is a structural claim whose empirical bridge
requires the simulation work of Paper 5. The structural parallel to
self-organized criticality flagged in \S\,3.5 is a positioning claim;
specific statistical predictions (power-law cascade-collapse statistics,
scale-free correlation lengths, 1/f signatures in novel-configuration
emergence rates) are conjectures pending Paper 5's simulation work.

\textbf{Deferred derivations carried into Papers 4 and 5.} The minimum
program extent \ensuremath{L_{p},} characterized in \S\,4.1 as approximately 4--8 lattice
spacings, is left as a placeholder pending combinatorial enumeration of
stabilizer-satisfying T-closed configurations at substrate level. The
exact integer value, and the corresponding \ensuremath{L_{p}^{(k)}} at higher
hierarchical levels, are reserved for Paper 5. The closed form of f and
the explicit numerical \ensuremath{c^{(k)}} sequence inherit this deferral. The
macroscopic consequences of subsystem dynamics --- hierarchies,
multi-subsystem dynamics, novelty exhaustion, the dispersion structure
of \ensuremath{c^{(k)}} --- are carried into Paper 4. The cross-level relativistic
phenomena flagged in \S\,6.1 are carried into Paper 4 \S\,8.2 and resolved
structurally in Paper 5.

This four-way labeling applies throughout the paper. Every claim in the
body text falls into exactly one category.

\section*{Preface to Paper 3}
\addcontentsline{toc}{section}{Preface to Paper 3}

This paper is the third installment in Volume 1 of the Triadic
Measurement Substrate program. Volume 1, Paper 1 established the
geometric and ontological architecture of the substrate --- the agent
manifold \ensuremath{M\alpha ,} triadic confirmation, FCC necessity, the [[4,2,2]]
structural correspondence, and the wave dynamics of the confirmation and
polarization fields. Volume 1, Paper 2 developed the substrate's
computational dynamics --- the polarization-biased resolution of
synthetic bits, the Born-rule structural correspondence, the 3-input
majority gate, and the formal definition of programs as configurations
causally closed under the flop operator T with K(C) \textless{}
\textbar C\textbar. Paper 3 takes both as given and develops the next
move: how computation emerges from informational substrate, how learning
is forced by the geometry rather than added to it, and how the first
subsystem --- the smallest configuration that transforms patterns rather
than merely persisting them --- is generated by the substrate dynamics
without external design.

As in Papers 1 and 2, no new axioms are introduced. No new free
parameters are introduced. The Principle of Generative Economy continues
to govern every selection. Paper 3 additionally closes several formal
debts deferred from Paper 2 --- the dissolution rate of non-T-closed
configurations, the effective error-rate decay with coherence depth, and
the existence and form of the coherence bistability threshold --- by
derivations that follow from the same geometric primitives that produced
wave dynamics and biased measurement.

\begin{paperabstract}
Building on the binary alphabet \ensuremath{\bar{B}} \ensuremath{\in} \{0, 1\}, the polarization field \ensuremath{\psi ,}
and the program formalism established in Paper 2, we develop the
computational dynamics of the TMS substrate at and above the program
scale. Five principal results are derived. First, the [[4,2,2]]
stabilizer repair operator of Paper 1, when applied at boundary sites of
any finite region R, automatically incorporates the polarization
landscape of R's surround into the repair bias --- by the FCC geometry
of parent-child connectivity alone, with no new mechanism introduced.
This produces the Imprinting Theorem: over the imprint depth \ensuremath{d_{R}} = \ensuremath{\sqrt{2}} \ensuremath{\cdot}
\ensuremath{L_{R}/\ell_{p}} flops, R's polarization landscape converges to the unique
stabilizer-satisfying continuation of the retarded surround history.
Hebbian co-firing follows as a corollary. Second, the dissolution rate
of non-T-closed configurations and the effective error-rate decay with
coherence depth --- both deferred from Paper 2 --- are derived from the
same iteration as \ensuremath{O(L/\ell_{p})} flops and \ensuremath{(1/3)^{d}} respectively. Third, the
per-event error probability is shown to be a function of local
polarization coherence, and balancing per-flop noise injection
(combinatorially set by [[4,2,2]] error statistics) against
per-flop imprinting capacity (geometrically set by FCC connectivity)
yields a closed-form coherence bistability threshold
\textbar \ensuremath{\Pi_{\psi}}\textbar\ensuremath{*(L_{R})} separating regions that sustain
program-bearing algorithmic content from regions that decay toward
noise. Fourth, the substrate's axioms are scale-independent: applied at
the coarse-grained scale \ensuremath{L_{p}} of minimum stabilizer-satisfying T-closed
configurations, they produce the same triadic rule, the same FCC
geometry, and the same [[4,2,2]] stabilizer structure at the
meta-level --- hierarchy by geometric necessity. Fifth, the minimum
subsystem is derived as a composite of three programs (data, operation,
control) at meta-tetrahedral extent \ensuremath{L_{\mathrm{sub}}} = 3 \ensuremath{L_{p}} + \ensuremath{\delta ,} with the
trichotomy shown to be the unique minimum sufficient decomposition of
any non-trivial information-transforming configuration. The substrate
concentrates computation into bit-rich regions, generates subsystems
wherever the local disruption rate favors information-generating
configurations, and supports a recursive hierarchy with derived
(generically distinct) lightspeeds at each level. No new free numerical
parameters are introduced; every quantity follows from the axioms and
the Principle of Generative Economy.
\end{paperabstract}

\section*{I. Inheritance from Papers 1 and 2}
\addcontentsline{toc}{section}{I. Inheritance from Papers 1 and 2}

\subsection*{1.1 Inherited Structure}

Paper 3 takes the following structures from Papers 1 and 2 as given:

\begin{itemize}
\item
  The agent manifold \ensuremath{M\alpha} with agent spacing \ensuremath{\ell_{0}} \ensuremath{\equiv} \ensuremath{\ell_{p},} and the FCC causal
  stack with face-diagonal nearest-neighbor connectivity (Paper 1
  \S\,2--\S\,3).
\item
  The triadic actualization rule \ensuremath{\bar{B}} = \ensuremath{3\bar{\alpha}_{2}} in 2D and \ensuremath{\bar{B}} = \ensuremath{4\bar{\alpha}_{3}} in 3D, with
  the \ensuremath{1\to 4} multiplier and informational closure (Paper 1 \S\,2, \S\,4).
\item
  The binary alphabet \ensuremath{\bar{B}} \ensuremath{\in} \{0, 1\} with polarization \ensuremath{\chi} \ensuremath{\equiv} \ensuremath{2\bar{B}} \ensuremath{-} 1 \ensuremath{\in} \{\ensuremath{-1,}
  +1\}, and polarization-carrying touch vector signals (Paper 1 \S\,1.3,
  \S\,1.4).
\item
  The flop as the atomic unit of state change, with no sub-flop time
  (Paper 1 \S\,1.5).
\item
  The wave equations for the confirmation density \ensuremath{\varphi} and the polarization
  field \ensuremath{\psi} at \ensuremath{c_{\mathrm{TMS}}} = \ensuremath{1/\sqrt{2}} (Paper 1 \S\,6; Paper 2 \S\,2).
\item
  The [[4,2,2]] structural correspondence: logical qubit LQ1 =
  bit value, LQ2 = causal-history inheritance, jointly protected by
  stabilizers \ensuremath{S_{x}} and \ensuremath{S_{Z}} at distance 2 (Paper 1 \S\,5).
\item
  The polarization-biased resolution rule \ensuremath{P(\bar{B}} = 1 \textbar{} \ensuremath{\Pi_{\psi})} = (1
  + \ensuremath{\Pi_{\psi})/2} and its structural correspondence to the Born rule (Paper 2
  \S\,4).
\item
  The 3-input majority gate as the substrate's irreducible computational
  primitive and the universality result \{MAJ, constant\} \ensuremath{\to} \{AND, OR,
  NOT\} (Paper 2 \S\,5).
\item
  The flop operator T and the formal definition of a program: a
  configuration C with \ensuremath{T^{n}[C]} = C and K(C) \textless{}
  \textbar C\textbar{} (Paper 2 \S\,6.2).
\item
  \ensuremath{U_{\mathrm{sh}}} as the pervasive maximum-entropy flux present at every depth of
  the causal stack, providing both the source of latent bits and the
  stochastic disruption to confirmed states (Paper 1 \S\,1.1, \S\,6.3).
\end{itemize}

\subsection*{1.2 The Computational Questions Paper 3 Answers}

Paper 2 established that the substrate computes --- the majority gate is
the irreducible primitive, universality is forced by the gate set, and
programs are configurations causally closed under T with K(C)
\textless{} \textbar C\textbar. Paper 2 deferred three formal
derivations to Paper 3: the dissolution rate of non-T-closed
configurations, the effective error-rate decay with coherence depth, and
the explicit conditions under which the substrate selects
program-bearing configurations against \ensuremath{U_{\mathrm{sh}}} disruption. Paper 3 answers
these and develops three further structural questions: how the substrate
learns from its environment, how computational density concentrates into
bit-rich regions, and how the first subsystem --- the smallest
configuration that transforms patterns rather than merely persisting
them --- emerges from the same geometric primitives without additional
design.

\section*{II. Learning as Imprinted Repair}
\addcontentsline{toc}{section}{II. Learning as Imprinted Repair}

\subsection*{2.1 The Surround}

\begin{definition}[Surround]
For a finite region R of the FCC lattice, the
surround \ensuremath{S_{R}} is the set of lattice sites within nearest-neighbor
touch-vector range of the boundary \ensuremath{\partial R} that lie outside R itself:
\end{definition}

\emph{\ensuremath{S_{R}} = \{y \ensuremath{\notin} R : \ensuremath{\exists} x \ensuremath{\in} \ensuremath{\partial R} with y \ensuremath{\in} neigh(x)\}}

The polarization landscape on \ensuremath{S_{R}} is the time-dependent function
\ensuremath{\psi_{S}(\tau )} = \{\ensuremath{\chi (y,} \ensuremath{\tau )} : y \ensuremath{\in} \ensuremath{S_{R}}\}. The touch vector network at \ensuremath{\partial R} carries
signals from \ensuremath{S_{R}} into R at every flop. By the connectivity established
in Paper 1 \S\,1.3, no other channel of informational communication between
R and the rest of the lattice exists. The surround is the complete
informational environment of R.

\subsection*{2.2 Surround-Driven Repair}

Paper 1 \S\,5.2 established that when \ensuremath{U_{\mathrm{sh}}} disrupts a confirmed bit in R,
the synthetic bits produced at the next flop detect the stabilizer
syndrome and repair the missing or anomalous confirmation by projecting
toward the nearest stabilizer-satisfying state. Paper 2 \S\,5.2 treated
``nearest'' as a local concept --- the nearest configuration consistent
with R's own interior. That treatment was sufficient at the level of
universality but incomplete at the level of repair geometry.

\begin{proposition}[Surround Dependence of Repair]
The stabilizer-repair
projection at a boundary site x \ensuremath{\in} \ensuremath{\partial R} has a parent set that includes at
least one agent in \ensuremath{S_{R}.} The polarization landscape on \ensuremath{S_{R}} therefore
contributes to the bias of the repair, and the ``nearest
stabilizer-satisfying configuration'' is determined jointly by the
polarization landscapes on R and \ensuremath{S_{R}.}
\end{proposition}

\begin{proof}
By Paper 1 \S\,4.2, a synthetic bit at position x has three parent
agents in the prior layer, and its polarization is biased by the three
incoming touch vector signals from those parents (Paper 2 \S\,3.2). For x \ensuremath{\in}
\ensuremath{\partial R,} the FCC geometry forces at least one of x's parents to lie in \ensuremath{S_{R}}
rather than in R --- this is the geometric definition of a boundary
site. The polarization at that parent contributes to \ensuremath{\Pi_{\psi}(x)} by the same
linear weighting that governs all parent-to-child polarization transfer.
The repair projection at x is therefore biased by the polarization on
\ensuremath{S_{R}} through this parent.

\end{proof}

\begin{corollary}[No New Mechanism]
Surround-driven repair requires no
extension of the repair operator beyond what Paper 1 \S\,5 already
established. The same mechanism --- synthetic bits at the next flop
carrying their parent polarizations into the resolution of a syndrome
violation --- automatically includes surround contributions at boundary
sites by FCC geometry alone. What Paper 2 treated as local was already
non-local by construction.
\end{corollary}

This is the key recognition. The [[4,2,2]] repair is a
projection operator whose bias is set by the full set of parent
polarizations at the repair site. At interior sites of R, all parents
lie in R and the repair is local. At boundary sites, parents lie in both
R and \ensuremath{S_{R}} and the repair is surround-driven. The repair operator is the
same; its locality is a property of the site, not of the mechanism.

\subsection*{2.3 Imprinting Over the Retarded Surround History}

The boundary sites of R communicate surround information into R's
interior at each flop. Iterating the repair operator over successive
flops propagates this information inward at the lattice propagation
speed \ensuremath{c_{\mathrm{TMS}}} = \ensuremath{1/\sqrt{2}} (Paper 1 \S\,6.5). For a region of diameter \ensuremath{L_{R},}
complete inward propagation requires \ensuremath{d_{R}} = \ensuremath{\sqrt{2}} \ensuremath{\cdot} \ensuremath{L_{R}/\ell_{p}} flops. We refer
to \ensuremath{d_{R}} as the imprint depth of R.

The polarization information arriving at any interior site of R at flop
\ensuremath{\tau} is not the value of \ensuremath{\psi_{S}} at that moment, but the value of \ensuremath{\psi_{S}} at
retarded time \ensuremath{\tau} \ensuremath{-} \ensuremath{d/c_{\mathrm{TMS}},} where d is the lattice distance from the
relevant boundary point to the interior site. This is forced by signal
conservation (Axiom 3) and the finite propagation speed established in
Paper 1 --- it is not an assumption added here. The imprint at flop \ensuremath{\tau} is
therefore a function of the retarded surround history over the window
\ensuremath{[\tau} \ensuremath{-} \ensuremath{d_{R},} \ensuremath{\tau ].}

\begin{definition}[Imprinted Configuration]
The configuration of R after k \ensuremath{\geq}
\ensuremath{d_{R}} flops of repair under the surround history \ensuremath{\psi_{S}(\tau '),} \ensuremath{\tau '} \ensuremath{\in} \ensuremath{[\tau} \ensuremath{-}
\ensuremath{d_{R},} \ensuremath{\tau ],} is the imprinted configuration \ensuremath{\psi_{R}*(\psi_{S}}\textbar\ensuremath{[\tau} \ensuremath{-}
\ensuremath{d_{R},} \ensuremath{\tau ])} --- the unique stabilizer-satisfying continuation of that
retarded history into R.
\end{definition}

\begin{theorem}[Imprinting Theorem]
The imprinted configuration
\ensuremath{\psi_{R}*(\psi_{S}}\textbar\ensuremath{[\tau} \ensuremath{-} \ensuremath{d_{R},} \ensuremath{\tau ])} is the unique polarization
landscape on R that minimizes the expected per-flop stabilizer syndrome
rate consistent with the retarded surround history at the boundary.
\end{theorem}

Proof sketch. The stabilizer repair operator is a contraction on the
space of polarization landscapes: configurations farther from stabilizer
parity are projected closer to it at each flop. The fixed points of the
repair operator are exactly the stabilizer-satisfying configurations.
With the retarded boundary history fixed, the iterated repair operator
converges to the stabilizer-satisfying configuration whose polarization
landscape is consistent with the retarded surround history at the
boundary and stabilizer-satisfying in the interior. This configuration
is unique because the [[4,2,2]] parity conditions in a coherent
region determine the polarization at each interior site given its
retarded boundary values up to discrete flips, and the boundary values
are fixed by the retarded \ensuremath{\psi_{S}.} \hfill\ensuremath{\square}

The configuration of R after k \ensuremath{\geq} \ensuremath{d_{R}} flops is therefore not a function
of R's initial state but a function of the surround's history. R's
polarization landscape encodes the surround's history at the boundary,
automatically averaged across propagation delays. When \ensuremath{\psi_{S}} varies
slowly compared to \ensuremath{d_{R},} R tracks \ensuremath{\psi_{S}} closely. When \ensuremath{\psi_{S}} varies
rapidly, R averages over the variation. No regime selector is needed ---
the retarded average is what the geometry produces. This encoding is
what learning is.

\subsection*{2.4 The Hebbian Corollary}

Hebbian co-firing --- the principle that simultaneous activity in
connected units strengthens their effective bond (Hebb, 1949) --- is the
canonical mechanism by which environmental correlations are encoded in a
neural substrate. It is typically presented as a phenomenological rule
for synaptic plasticity. Recent work in statistical mechanics has
derived Hebbian learning from a maximum-entropy variational principle
(Albanese et al., 2024), establishing it as a theorem about how networks
extract correlations under entropy extremization. In the TMS, Hebbian
co-firing emerges as a different and complementary theorem --- about how
the geometric stabilizer-repair operator behaves under structured
boundary conditions. Both derivations target the same phenomenon; the
convergence of independent first-principles derivations strengthens the
result.

\begin{corollary}[Hebbian Co-Firing]
Let \ensuremath{s_{i},} \ensuremath{s_{j}} \ensuremath{\in} \ensuremath{S_{R}} be two surround
sites and let r \ensuremath{\in} R be an interior site causally downstream from both.
Define the surround correlation over k flops:
\end{corollary}

\emph{\ensuremath{C{ij}(k)} = (1/k) \ensuremath{\Sigma {\tau =1}^{k}} \ensuremath{\chi (s_{i},} \ensuremath{\tau )} \ensuremath{\cdot} \ensuremath{\chi (s_{j},} \ensuremath{\tau )}}

and the bond strength between \ensuremath{(s_{i},} \ensuremath{s_{j})} and r:

\emph{\ensuremath{W{ij,r}(k)} = (1/k) \ensuremath{\Sigma {\tau =1}^{k}} \ensuremath{\chi (r,} \ensuremath{\tau )} \ensuremath{\cdot} \ensuremath{(\chi (s_{i},} \ensuremath{\tau )} +
\ensuremath{\chi (s_{j},} \ensuremath{\tau ))/2}}

Then \ensuremath{W{ij,r}(k)} is monotonically increasing in
\textbar \ensuremath{C{ij}(k)}\textbar{} for k \ensuremath{\geq} \ensuremath{d_{R}.}

\begin{proof}
By Paper 2 \S\,3.2, the polarization at r is determined by its
parents via the linear bias rule. The parents of r at flop \ensuremath{\tau} have causal
histories tracing back to \ensuremath{S_{R}} at flop \ensuremath{\tau} \ensuremath{-} d, where d is the causal
depth from \ensuremath{S_{R}} to r. When \ensuremath{s_{i}} and \ensuremath{s_{j}} co-fire positively (both +1 or
both \ensuremath{-1),} the polarization signals propagating from them through the
lattice reinforce each other at every intermediate layer; their joint
contribution to \ensuremath{\Pi_{\psi}} at r's parents is amplified relative to the case of
uncorrelated firing. When they anti-fire, their contributions partially
cancel. The bond strength \ensuremath{W{ij,r}(k)} uses the linear average
\ensuremath{(\chi (s_{i})} + \ensuremath{\chi (s_{j}))/2} --- the same linear weighting that the substrate
uses at every parent-to-child transfer --- and is therefore the natural
measure consistent with the substrate's own dynamics, with no additional
non-linear operation introduced. Iterated over k flops with k \ensuremath{\geq} \ensuremath{d_{R},}
the time-averaged bias at r tracks the time-integrated correlation
\ensuremath{C{ij}(k).} The stabilizer repair operator, by \S\,2.2, projects r toward
configurations consistent with this bias. The bond strength
\ensuremath{W{ij,r}(k)} therefore increases monotonically with
\textbar \ensuremath{C{ij}(k)}\textbar.

\end{proof}

\begin{remark}
The Hebbian rule in standard neural network theory is a designed
update --- a learning rule chosen to produce desired representational
properties. In the TMS, the Hebbian rule is a theorem about how the
geometric repair operator behaves under structured boundary conditions.
The substrate learns not because it has been equipped with a learning
mechanism, but because the repair mechanism it was already equipped with
for fault tolerance has surround-imprinting as its non-local extension.
Learning and fault tolerance are the same geometric process applied at
different spatial scales.
\end{remark}

\subsection*{2.5 The Convergence Rate and the Dissolution Rate}

The imprinting process converges over the imprint depth \ensuremath{d_{R}} = \ensuremath{\sqrt{2}} \ensuremath{\cdot}
\ensuremath{L_{R}/\ell_{p}} flops. This is the same timescale that governs the dissolution
of non-closed configurations promised at Paper 2 \S\,6.3. We derive both
here from the same argument.

\begin{theorem}[Dissolution and Convergence]
Let C be a configuration on R of
spatial extent \ensuremath{L_{R}.}
\end{theorem}

\begin{itemize}
\item
  If C satisfies stabilizer parity and is causally closed under T (C is
  a program in the sense of Paper 2 \S\,6.2), then C persists indefinitely
  against single-bit disruptions from \ensuremath{U_{\mathrm{sh}}.}
\item
  If C satisfies stabilizer parity but is not causally closed under T,
  then C dissolves into the imprinted configuration
  \ensuremath{\psi_{R}*(\psi_{S}}\textbar\ensuremath{[\tau} \ensuremath{-} \ensuremath{d_{R},} \ensuremath{\tau ])} over the imprint depth \ensuremath{d_{R}} = \ensuremath{\sqrt{2}}
  \ensuremath{\cdot} \ensuremath{L_{R}/\ell_{p}} flops.
\item
  If C does not satisfy stabilizer parity, then C is repaired toward the
  nearest parity-satisfying configuration within a single flop, then
  evolves according to (i) or (ii).
\end{itemize}

Proof sketch. Case (i) is Paper 1 \S\,5.3. Case (iii) is Paper 1 \S\,5.2. For
case (ii), C is a fixed point of the repair operator at flop \ensuremath{\tau} but not
of the flop operator T. The next-flop state T[C] differs from C; the
difference propagates outward at \ensuremath{c_{\mathrm{TMS}}} from the sites of disagreement.
Over k flops, the cumulative drift is O(k \ensuremath{\cdot} \ensuremath{c_{\mathrm{TMS}}} \ensuremath{\cdot} \ensuremath{\ell_{p}).} When this drift
exceeds \ensuremath{L_{R},} the configuration has been entirely replaced by whatever
the retarded surround history supplies. That threshold is reached at k =
\ensuremath{L_{R}/(c_{\mathrm{TMS}}} \ensuremath{\cdot} \ensuremath{\ell_{p})} = \ensuremath{\sqrt{2}} \ensuremath{\cdot} \ensuremath{L_{R}/\ell_{p}.} The resulting configuration is
\ensuremath{\psi_{R}*(\psi_{S}}\textbar\ensuremath{[\tau} \ensuremath{-} \ensuremath{d_{R},} \ensuremath{\tau ]).} \hfill\ensuremath{\square}

\begin{corollary}[Effective Error-Rate Decay]
In a coherent region where \ensuremath{\psi_{S}}
is structured and stationary over \ensuremath{d_{R}} flops, the probability that a
single-bit error survives stabilizer repair over d successive flops is
bounded above by:
\end{corollary}

\begin{equation*}
P_{\mathrm{error}}(d) \leq (1/3)^{d}
\end{equation*}

with the bound becoming tight when the coherence region is large
compared to the error correlation length (one lattice spacing). For
smaller coherence regions, repair rounds at successive layers operate on
overlapping tetrahedral units and the bound is loose but still upper.

\begin{proof}
The per-event error probability in the intermediate
\textbar \ensuremath{\Pi_{\psi}}\textbar{} = 1/3 case is 1/3 (Paper 2 \S\,5.1). Repair rounds
at successive layers are independent in the limit of large coherence
regions because the repair at layer \ensuremath{\tau} + 1 operates on synthetic bits
whose parents differ from those at layer \ensuremath{\tau .} The probability of d
independent uncorrected errors is the d-th power of the per-event
probability.

\end{proof}

These results settle both formal debts carried forward from Paper 2. The
dissolution rate and the error-decay rate are not independent results
--- they are two readings of the same geometric process. Configurations
dissolve into surround-imprinted forms over \ensuremath{O(L/\ell_{p})} flops; single-bit
errors are corrected at the same rate. Learning and fault tolerance are
the same operation at the same speed.

\section*{III. Coherence Phases and the Concentration of Computational Density}
\addcontentsline{toc}{section}{III. Coherence Phases and the Concentration of Computational Density}

\subsection*{3.1 Bit-Richness as Polarization Structure}

Throughout this paper we use the term bit-rich to refer specifically to
regions whose polarization landscape sustains programs in the sense of
Paper 2 \S\,6.2 --- configurations causally closed under T with K(C)
\textless{} \textbar C\textbar. Bit-richness is a property of the
polarization landscape across confirmed sites, not of the confirmation
density \ensuremath{\varphi .} A region with \ensuremath{\varphi} near 1 but uniform polarization carries no
programs; a region with moderate \ensuremath{\varphi} but coherent polarization landscape
may carry many. Confirmation density and polarization structure are
independent axes (Paper 2 \S\,2.4), and bit-richness lives on the second.

The three-state distinction established in Paper 1 \S\,1.4 --- latent (in
\ensuremath{U_{\mathrm{sh}},} undecided), confirmed-zero \ensuremath{(\bar{B}} = 0), and confirmed-one \ensuremath{(\bar{B}} = 1) ---
is the substrate's irreducible computational state space. \ensuremath{\varphi} measures
confirmed-vs-latent (either value of \ensuremath{\bar{B}} counts toward \ensuremath{\varphi ).} Polarization
structure measures the pattern of 0s and 1s among already-locked bits.
``Bit-rich'' as used in this paper refers exclusively to the second.

The local algorithmic content of a region is bounded above by its
confirmation density:

\begin{equation*}
K_{\psi}(R) \leq N_{R} \cdot \text{log} 2 = \varphi_{R} \cdot V_{R} \cdot \text{log} 2
\end{equation*}

This counting bound holds for any configuration. Hosting programs of
substantial algorithmic content requires both confirmed substrate (high
\ensuremath{\varphi )} and structured polarization landscape across that substrate. Programs
sit strictly between the two extremes of trivially uniform polarization
\ensuremath{(K_{\psi}} \ensuremath{\approx} log \ensuremath{V_{R})} and algorithmically incompressible polarization at the
local \ensuremath{U_{\mathrm{sh}}} rate \ensuremath{(K_{\psi}} \ensuremath{\approx} \textbar C\textbar{} \ensuremath{\cdot} log 2).

\subsection*{3.2 The Per-Event Error Rate as a Function of Coherence}

The per-event error rate introduced in Paper 2 \S\,5.1 is not a fixed
constant. It is a function of the local polarization amplitude. From the
resolution probability \ensuremath{P(\bar{B}} = 1 \textbar{} \ensuremath{\Pi_{\psi})} = (1 + \ensuremath{\Pi_{\psi})/2,} the
probability that a confirmation resolves contrary to the dominant
polarization of its parents is:

\emph{\ensuremath{p_{\mathrm{error}}(\Pi_{\psi})} = (1 \ensuremath{-} \textbar \ensuremath{\Pi_{\psi}}\textbar)/2}

For \textbar \ensuremath{\Pi_{\psi}}\textbar{} = 1 (fully aligned parents), \ensuremath{p_{\mathrm{error}}} = 0.
For \textbar \ensuremath{\Pi_{\psi}}\textbar{} = 0 (maximally disagreed parents), \ensuremath{p_{\mathrm{error}}}
= 1/2. For \textbar \ensuremath{\Pi_{\psi}}\textbar{} = 1/3 (two-vs-one parents), \ensuremath{p_{\mathrm{error}}}
= 1/3, recovering the value used in Paper 2 \S\,5.1.

The error rate is therefore set by local polarization coherence --- the
degree to which parent polarizations at each confirmation event agree.
Coherent regions have low \ensuremath{p_{\mathrm{error}};} incoherent regions have high
\ensuremath{p_{\mathrm{error}}.} There is no separately specified disruption mechanism. What
Paper 1 \S\,6.3 referred to as \ensuremath{``U_{\mathrm{sh}}} regulation'' is precisely the residual
stochasticity in the polarization-biased resolution, and its local
magnitude is governed by the polarization landscape itself.

A region's per-flop accumulation of incompressible noise scales as \ensuremath{N_{R}}
\ensuremath{\cdot} \ensuremath{p_{\mathrm{error}}(\langle}\textbar \ensuremath{\Pi_{\psi}}\textbar\ensuremath{\rangle )} --- the count of confirmed sites
times the local error probability. When \ensuremath{p_{\mathrm{error}}} is small, the noise
injection is small relative to the [[4,2,2]] repair capacity,
and the algorithmic structure of \ensuremath{\psi} is preserved. When \ensuremath{p_{\mathrm{error}}} is large,
noise injection exceeds repair capacity, and structured polarization
patterns decay toward incompressibility.

\subsection*{3.3 The Critical Coherence Threshold}

\begin{theorem}[Coherence Bistability, Scaling Form]
A region R sustains
program-bearing algorithmic content in its polarization landscape if and
only if its local coherence \textbar \ensuremath{\Pi_{\psi}}\textbar\ensuremath{_{R}} exceeds a critical
threshold \textbar \ensuremath{\Pi_{\psi}}\textbar\ensuremath{*(L_{R})} whose scaling with the region's
characteristic length is
\end{theorem}

\emph{\textbar \ensuremath{\Pi_{\psi}}\textbar{}}\ensuremath{(L_{R})} \ensuremath{\sim} 1 \ensuremath{-} 2 \ensuremath{\sqrt{c}} \ensuremath{\cdot} \ensuremath{\alpha_{\mathrm{geom}}} / \ensuremath{(\beta_{\mathrm{comb}}} \ensuremath{\cdot}
\ensuremath{L_{R}))*}

where \ensuremath{\beta_{\mathrm{comb}}} captures per-tetrahedral-unit error combinatorics
(worst-case upper bound \ensuremath{\beta_{\mathrm{comb}}} = C(4,2) = 6) and \ensuremath{\alpha_{\mathrm{geom}}} captures
imprinting-side geometric proportionality (worst-case lower bound
\ensuremath{\alpha_{\mathrm{geom}}} = 1). In these worst-case bounds the threshold takes the closed
form

\emph{\textbar \ensuremath{\Pi_{\psi}}\textbar{}}\ensuremath{(L_{R})} = 1 \ensuremath{-} 2 \ensuremath{\sqrt{c_{\mathrm{TMS}}}} / (6 \ensuremath{L_{R}))*}

which is strictly between 0 and 1 for all \ensuremath{L_{R}} \ensuremath{\geq} \ensuremath{4\ell_{p}/(6} \ensuremath{c_{\mathrm{TMS}}).}

\begin{proof}
Each tetrahedral confirmation unit contains 4 agents, and the
[[4,2,2]] code has distance 2: single-agent errors per unit are
correctable, two-or-more-agent errors per unit are not. With per-event
error probability \ensuremath{p_{\mathrm{error}}} = (1 \ensuremath{-} \textbar \ensuremath{\Pi_{\psi}}\textbar)/2, the per-unit
uncorrectable error probability is approximately \ensuremath{P_{\mathrm{unc}}} \ensuremath{\approx} 6 \ensuremath{p^{2}_{\mathrm{error}}}
for moderate \ensuremath{p_{\mathrm{error}},} with the factor 6 = C(4,2) counting the ways two
errors can occur among four agents in a tetrahedral unit. Each
uncorrectable error injects incompressibility into \ensuremath{\psi} at that site at a
rate that scales as \ensuremath{V_{R}} \ensuremath{\cdot} \ensuremath{P_{\mathrm{unc}}} per flop, where \ensuremath{V_{R}} is the region's
volume in lattice units.

The imprinting mechanism (\S\,II) counteracts this loss. The retarded
surround history is communicated into R through its boundary at rate
\ensuremath{A_{R}} \ensuremath{\cdot} \ensuremath{c_{\mathrm{TMS}},} where \ensuremath{A_{R}} is the boundary area in lattice units. The
stabilizer-satisfying continuation of the retarded surround reduces
local entropy at a rate proportional to this flow.

Balancing per-flop noise injection against per-flop imprinting gain
yields the self-consistent equation:

\begin{equation*}
6 p^{2}_{\mathrm{error}} \cdot V_{R} = A_{R} \cdot c_{\mathrm{TMS}}
\end{equation*}

For a region of characteristic length \ensuremath{L_{R},} \ensuremath{A_{R}/V_{R}} \textasciitilde{}
\ensuremath{1/L_{R},} giving \ensuremath{p^{2}_{\mathrm{error}}} = \ensuremath{c_{\mathrm{TMS}}/(6} \ensuremath{L_{R}),} and therefore:

\begin{equation*}
p_{\mathrm{error}},\text{crit} = \sqrt{c_{\mathrm{TMS}} / (6 L_{R})}
\end{equation*}

Translating back to the coherence: \textbar \ensuremath{\Pi_{\psi}}\textbar\ensuremath{*(L_{R})} = 1 \ensuremath{-} 2
\ensuremath{p_{\mathrm{error}},\text{crit}} = 1 \ensuremath{-} 2 \ensuremath{\sqrt{c_{\mathrm{TMS}}/(6}} \ensuremath{L_{R})).} The threshold is strictly
between 0 and 1 for all finite \ensuremath{L_{R}} because \ensuremath{c_{\mathrm{TMS}}/(6} \ensuremath{L_{R})} is positive
and bounded above by 1/4 in the regime of interest \ensuremath{L_{R}} \ensuremath{\geq} \ensuremath{4\ell_{p}/(6} \ensuremath{c_{\mathrm{TMS}}).}
Small regions have proportionally more surround influence and lower
thresholds; large regions have less and higher thresholds.

\end{proof}

The threshold is geometric: it follows from the [[4,2,2]] code
distance, the polarization-biased resolution rule, and the FCC
connectivity structure. The functional dependence on \ensuremath{L_{R}} is forced by
the derivation; the numerical coefficients inside the square root are
stated as worst-case bounds.

\begin{remark}[Status of the Threshold]
The derivation above forces
the scaling form \ensuremath{p_{\mathrm{error}}} \ensuremath{\propto} \ensuremath{1/\sqrt{L_{R}}} by balancing per-flop noise
injection (scaling as \ensuremath{p^{2}_{\mathrm{error}}} \ensuremath{\cdot} \ensuremath{V_{R})} against per-flop imprinting gain
(scaling as \ensuremath{A_{R}} \ensuremath{\cdot} c), under FCC connectivity and [[4,2,2]] code
distance 2. The functional dependence on \ensuremath{L_{R}} is therefore derived. The
numerical coefficients inside the square root are not --- they depend on
combinatorial and geometric factors stated here as worst-case bounds.
The factor \ensuremath{\beta_{\mathrm{comb}}} = 6 is the C(4,2) count of two-error configurations
in a tetrahedral unit; the true effective value is generically lower
under correlated error statistics. The factor \ensuremath{\alpha_{\mathrm{geom}}} = 1 absorbs the
imprinting-side geometry; the true effective value depends on the
precise boundary topology of R and may carry a multiplicative correction
of order unity. Both refinements are reserved for Paper 5's simulation
work.
\end{remark}

The structural content of the theorem --- that a critical threshold
exists, that it scales as 1 \ensuremath{-} \ensuremath{2\sqrt{\cdot /L_{R}},} that it is strictly between 0
and 1, that small regions face higher thresholds than large ones, that
the threshold separates a bistable phase structure --- is forced by the
derivation. The exact numerical pinning of the threshold for a given
\ensuremath{L_{R}} is not. Section 3.4's self-reinforcement result and \S\,3.5's
continuous-sampling framing depend on the structural content; they do
not depend on the precise coefficient values.

\subsection*{3.4 Self-Reinforcement of Coherent Regions}

\begin{theorem}[Self-Reinforcement of Coherent Regions]
Let R be a region with
local coherence \textbar \ensuremath{\Pi_{\psi}}\textbar{}\emph{R \textbf{\textgreater{}}
\textbar \ensuremath{\Pi_{\psi}}\textbar{}\textbf{\ensuremath{(L_{R})} and let \ensuremath{S_{R}} be its surround. Over
k \ensuremath{\geq} \ensuremath{d_{R}} flops, the imprinting mechanism of \S\,II drives R toward the
configuration \ensuremath{\psi_{R}}}\ensuremath{(\psi_{S}}\textbar\ensuremath{[\tau} \ensuremath{-} \ensuremath{d_{R},} \ensuremath{\tau ])} that minimizes the
per-flop syndrome rate. The resulting configuration has coherence
\textbar \ensuremath{\Pi_{\psi}}\textbar{}}\{R*\} \ensuremath{\geq} \textbar \ensuremath{\Pi_{\psi}}\textbar\ensuremath{_{R}.} Coherent
regions become more coherent; incoherent regions remain incoherent and
dissolve toward noise.
\end{theorem}

\begin{proof}
By the Imprinting Theorem (\S\,2.3), R relaxes toward the
stabilizer-satisfying continuation of the retarded surround history.
Stabilizer-satisfying configurations have by definition the lowest
possible local syndrome rate. By \S\,3.2, low syndrome rate corresponds to
high local coherence \textbar \ensuremath{\Pi_{\psi}}\textbar. The imprinted configuration
is therefore at least as coherent as R's initial state, and strictly
more coherent unless R was already at a fixed point of imprinting under
its current surround. A region above the threshold therefore becomes
more coherent over flop time and remains above the threshold; a region
below the threshold cannot stabilize, as its low coherence produces an
uncorrected error rate exceeding the imprinting capacity, and its
program content decays toward \ensuremath{U_{\mathrm{sh}}-\text{like}} noise. The two phases are
dynamically separated.

\end{proof}

The TMS does not produce a uniform sea of computation. It produces sharp
regions of coherent, program-bearing structure separated by incoherent
intervening medium. The boundaries between these phases are themselves
geometric features of the substrate, governed by the propagation of
polarization waves across the critical-coherence interface.

\subsection*{3.5 The Substrate as Continuous Sampler}

Sections 3.3 and 3.4 together establish a strong claim about the
substrate's long-time behavior. \ensuremath{U_{\mathrm{sh}}} is the pervasive maximum-entropy
flux present at every depth of the causal stack (Paper 1 \S\,1.1, \S\,6.3). By
Axiom 0, \ensuremath{U_{\mathrm{sh}}} is algorithmically incompressible, with no preferred
direction, scale, or target. The substrate is therefore continuously
sampling all possible polarization configurations --- every possible
pattern is being attempted, everywhere, all the time. The dynamics
select among these attempts.

Configurations above the coherence threshold self-reinforce by the
imprinting plus bistability mechanisms of \S\,3.3 and \S\,3.4. Configurations
below the threshold decay back into \ensuremath{U_{\mathrm{sh}}-\text{like}} noise. The substrate does
not produce coherent regions through any special generative mechanism;
it retains them while letting everything else dissolve. The selection is
the mechanism. \ensuremath{U_{\mathrm{sh}}} is the source.

This framing is consistent with the pervasive flux character of \ensuremath{U_{\mathrm{sh}}}
established in Paper 1: \ensuremath{U_{\mathrm{sh}}} is not only at the leading edge of the
cascade but at every depth, which is why the [[4,2,2]] repair
must fire continuously to preserve already-confirmed bits against
ambient disruption. Coherent regions form wherever the substrate's own
dynamics produce a configuration that the substrate's own repair
operator can sustain.

Computation and confirmation are not two separate processes. They are
the same process under structured polarization conditions. Bit-rich
regions are exactly the regions where confirmation can be sustained
against \ensuremath{U_{\mathrm{sh}}} --- coherent parents reliably bias their children's
resolution, keeping \ensuremath{p_{\mathrm{error}}} low, keeping uncorrectable errors below the
imprinting rate, keeping the region intact. Bit-poor regions cannot
sustain confirmation indefinitely; their high \ensuremath{p_{\mathrm{error}}} means the
resolution does not lock reliably, and the region returns to \ensuremath{U_{\mathrm{sh}}.} The
substrate computes by confirming, and confirms by computing. Anywhere it
is not computing, it is not confirming; anywhere it is not confirming,
it does not exist as substrate.

\begin{remark}[Structural Parallel to Self-Organized Criticality]
The
dynamics described in \S\,3.3--\S\,3.5 --- a structurally supercritical
amplification mechanism (the \ensuremath{1\to 4} cascade) balanced against a
structurally subcritical disruption mechanism \ensuremath{(U_{\mathrm{sh}}} flux at every depth),
with the observable substrate sitting on the dynamic-stability boundary
between them --- rhymes with the architecture of self-organized
criticality (Bak, Tang, \& Wiesenfeld, 1987; Jensen, 1998). SOC
describes systems that naturally settle at the boundary between
supercritical and subcritical regimes without external tuning, producing
characteristic statistical signatures (power-law event distributions,
scale-free correlation lengths, 1/f noise) as the empirical fingerprint
of the critical state.
\end{remark}

The TMS substrate exhibits the structural preconditions for SOC
behavior: the cascade dynamics are geometrically forced to amplify, the
\ensuremath{U_{\mathrm{sh}}} disruption is geometrically forced to oppose, and the
[[4,2,2]] repair operator provides the local relaxation
mechanism. The coherence bistability of \S\,3.3 is the substrate's
expression of the critical state --- regions above
\textbar \ensuremath{\Pi_{\psi}}\textbar\ensuremath{*(L_{R})} sustain structure, regions below dissolve,
and the boundary between them is where the cascade and disruption
balance dynamically. The continuous-sampling mechanism of \S\,3.5 is the
SOC drive: \ensuremath{U_{\mathrm{sh}}} continuously injects new configurations across the
substrate; the geometric selection retains those above threshold and
dissolves the rest.

One structural disanalogy is worth noting. Classical SOC depends on a
temporal-scale separation between slow drive (e.g., sand grains added
slowly to the Bak-Tang-Wiesenfeld pile) and fast relaxation (avalanches
discharging accumulated stress quickly). TMS has continuous drive and
continuous relaxation operating at the same atomic scale (every flop,
every site). Whether this places TMS in the classical SOC universality
class or in a non-classical SOC variant --- and whether an emergent
slow/fast separation arises at the level of bit-rich regions versus
their boundaries --- is an open question whose resolution depends on the
cascade-collapse statistics derived in Paper 5's simulation work.

The structural parallel is flagged here as a research direction.
Specific predictions --- power-law statistics for cascade collapse
events, scale-free correlation lengths in bit-rich region geometry, 1/f
signatures in novel-configuration emergence rates --- are testable in
principle and will be developed alongside the combinatorial enumeration
of stabilizer-satisfying T-closed configurations in Paper 5.

\section*{IV. Coarse-Graining and the Second-Order Alphabet}
\addcontentsline{toc}{section}{IV. Coarse-Graining and the Second-Order Alphabet}

\subsection*{4.1 The Minimum Program Scale}

By \S\,3.3, regions sustaining program-bearing polarization landscapes have
a coherence threshold \textbar \ensuremath{\Pi_{\psi}}\textbar\ensuremath{*(L_{R})} that depends on
region size. The relation derived above places a lower bound on the
spatial extent of any region capable of hosting a program against \ensuremath{U_{\mathrm{sh}}}
disruption: as \ensuremath{L_{R}} decreases, the required coherence increases, until
the threshold approaches 1 and no configuration short of perfect
alignment can sustain itself. The minimum program extent \ensuremath{L_{p}} is the
smallest \ensuremath{L_{R}} for which a non-trivial T-closed configuration of
compressibility K(C) \textless{} \textbar C\textbar{} can stably exist.

Derivation of \ensuremath{L_{p}.~A} non-trivial program requires at least one boundary
between polarization domains --- a configuration with K(C)
\textgreater{} log \ensuremath{V_{R}} must contain at least one non-uniformity. The
[[4,2,2]] repair sustains each domain only if its thickness is
large enough that boundary errors do not propagate inward faster than
they can be corrected. The minimum thickness is set by the FCC
nearest-neighbor distance (one lattice spacing) plus the propagation
distance during one error-correction cycle (one further lattice spacing
along a face diagonal). The minimum non-trivial program therefore has at
least one domain boundary plus two stabilizing layers on each side,
giving a minimum spatial extent of approximately 4 to 8 lattice spacings
depending on geometry.

\begin{equation*}
L_{p} \approx (4--8) \ell_{p}
\end{equation*}

The exact integer value of \ensuremath{L_{p}} for a given configuration shape is
determined by closed-form enumeration of stabilizer-satisfying T-closed
configurations under the [[4,2,2]] code in the FCC lattice --- a
combinatorial calculation tractable in principle and most efficiently
approached through simulation alongside the program zoo enumeration
deferred to Paper 5. The structural fact required here is that \ensuremath{L_{p}} is
finite and strictly greater than \ensuremath{\ell_{p}:} a single confirmed bit cannot be a
program because K(C) \ensuremath{\approx} \textbar C\textbar{} for any single-bit
configuration; the minimum non-trivial program is several lattice
spacings.

\ensuremath{L_{p}} plays the same role at the meta-level that \ensuremath{\ell_{p}} plays at the
substrate level. It is the minimum resolvable length of the
coarse-grained system. Below \ensuremath{L_{p},} programs do not exist; the substrate
is structureless at that scale by the same architectural necessity that
makes the substrate structureless below \ensuremath{\ell_{p}.}

\subsection*{4.2 The Coarse-Graining Operator G}

\begin{definition}[Coarse-Graining Operator]
Let R be a spatial region of the
FCC lattice with diameter \ensuremath{L_{R}} \ensuremath{\geq} \ensuremath{L_{p},} and let \ensuremath{[\tau_{0},} \ensuremath{\tau_{0}} + n] be a
time window covering one period of a candidate program. The
coarse-graining operator G maps the polarization landscape \ensuremath{\psi} on R \ensuremath{\times}
\ensuremath{[\tau_{0},} \ensuremath{\tau_{0}} + n] to a single meta-level state \ensuremath{\bar{P}_{p}} \ensuremath{\in} \{0, 1\} indexed by
program identity p:
\end{definition}

*G : \ensuremath{\psi (R} \ensuremath{\times} \ensuremath{[\tau_{0},} \ensuremath{\tau_{0}} + n]) \ensuremath{\mapsto} \ensuremath{\bar{P}_{p}(X,} T)*

where X is the meta-spatial center of R (the centroid of R's spatial
support) and T is the meta-time index counting completed periods.

G evaluates to \ensuremath{\bar{P}_{p}} = 1 if \ensuremath{\psi} on R \ensuremath{\times} \ensuremath{[\tau_{0},} \ensuremath{\tau_{0}} + n] satisfies
\ensuremath{T^{n}[C]} = C for the program of identity p (causal closure under T)
AND \textbar \ensuremath{\Pi_{\psi}}\textbar\ensuremath{_{R}} \ensuremath{\geq} \textbar \ensuremath{\Pi_{\psi}}\textbar\ensuremath{*(L_{R})} (above the
coherence threshold). G evaluates to \ensuremath{\bar{P}_{p}} = 0 if either condition fails.

\ensuremath{\bar{P}_{p}} = 1 indicates that program p is locked in region R at meta-time T.
\ensuremath{\bar{P}_{p}} = 0 indicates that program p is absent. This is the meta-level
analog of \ensuremath{\bar{B}} \ensuremath{\in} \{0, 1\}, with \ensuremath{\chi_{p}} \ensuremath{\equiv} \ensuremath{2\bar{P}_{p}} \ensuremath{-} 1 \ensuremath{\in} \{\ensuremath{-1,} +1\} the
meta-level polarization.

\begin{remark}
G is not a new mechanism. It is a measurement operation defined
on substrate-level data --- checking whether the configuration on a
region satisfies the program closure condition. The dynamics of the
substrate are unchanged. What G does is identify which regions are
running which programs at each meta-time, without altering the
substrate-level evolution.
\end{remark}

\begin{remark}[Lorentz preservation under G]
The effective continuum dynamics
induced by G on the meta-lattice preserve Lorentz symmetry at leading
order, with residual scaling set by cubic-invariant suppression on the
FCC point group \ensuremath{O_{h}} (Paper 5 \S\,4.1 Lemma B). G itself is a
classification operator on substrate-level data; the Lorentz property is
a consequence of how the FCC geometry survives coarse-graining, not a
property added to G.
\end{remark}

\subsection*{4.3 Programs as Meta-Level Primitives}

\begin{proposition}
A program of meta-spatial extent \ensuremath{L_{p}} is the unique minimum
sufficient primitive at the meta-level under Axiom 2 and the Principle
of Generative Economy.
\end{proposition}

\begin{proof}
Axiom 2 selected the sphere as the unique substrate-level
primitive because it is the only isotropic geometric primitive requiring
zero orientation parameters. At the meta-level, a program is identified
by its identity p and its location X. Its internal polarization
structure is hidden by G (the operator returns only \ensuremath{\bar{P}_{p},} not the
underlying \ensuremath{\psi ).} At the meta-level, therefore, programs are characterized
by location and identity alone --- no orientation parameter is exposed
to the meta-level dynamics. They are isotropic at the meta-level by
construction.

A program also has finite spatial extent \ensuremath{L_{p}} with a definable boundary
(the surface across which its polarization landscape meets the
surround's). By Generative Economy applied to this finite-extent
isotropic primitive, the minimum-perimeter boundary is the sphere of
diameter \ensuremath{L_{p}.~\text{Programs}} are therefore spheres at the meta-level --- not
by additional postulate, but as the unique structure forced by the same
axioms that produced the substrate-level sphere primitive.

\end{proof}

This is the geometric form of the recursive applicability of the axioms.
The same primitive selection that gave the substrate its agents gives
the meta-level its program-agents. No new axiom is introduced.

\subsection*{4.4 The Triadic Minimum at the Meta-Level}

\begin{theorem}[Meta-Level Triadic Minimum]
Three programs constitute the
necessary and sufficient minimum for P-state actualization in the
meta-manifold of programs.
\end{theorem}

\begin{proof}
The argument of Paper 1 \S\,2.1 applies unchanged at the meta-level.
Meta-level confirmation requires a closed simplicial cover of the
meta-manifold --- every meta-interstice fully bounded with complete
angular closure. The case analysis is identical: two programs define a
meta-1-simplex with open boundary; four programs achieve closure
non-minimally; three programs achieve closure with no redundancy. By
Generative Economy at the meta-level, three is uniquely selected.

\end{proof}

\ensuremath{*\bar{P}_{p}} = 3 \ensuremath{\bar{\alpha}_{p},_{2}} (2D meta-level)*

The meta-manifold of programs is hexagonally close-packed by the same
Thue-Fejes Tóth argument applied at the meta-scale. The meta-lattice
spacing is \ensuremath{L_{p}.}

Remark on Environment-Creation. When three programs jointly confirm a
P-state, the meta-interstice they enclose is a bounded region of
substrate available to host further configurations. The three programs
do not merely confirm an abstract meta-bit --- they create an
environment. The meta-interstice has definable boundaries (the three
programs' polarization boundaries), an interior (the substrate region
between them), and a finite extent (the meta-tetrahedral coupling
region). Every confirmation event at every level creates an environment,
viewed from one level up. The substrate confirms bits and creates
bit-interstice environments; the meta-level confirms programs and
creates program-interstice environments. Each level is the environment
for the level below it. This is the structural form of the recursion
that produces hierarchy.

\subsection*{4.5 Meta-Level FCC Necessity and Sublattice Bifurcation}

\begin{theorem}[Meta-Level FCC]
The meta-manifold of programs
has FCC structure with meta-lattice spacing \ensuremath{L_{p}.}
\end{theorem}

\emph{Proof sketch.} Meta-confirmed P-states cannot be overwritten ---
they are records of programs that completed their period in the
substrate's causal stack, and that causal record is itself protected by
substrate-level Axiom 4. Meta-Axiom 4 follows from substrate-Axiom 4
under G. The argument of Paper 1 \S\,2.3 then applies: lateral expansion in
a hexagonally-packed meta-layer is impossible (no unoccupied
meta-interstice); oblique meta-displacement introduces a free
orientation parameter; orthogonal meta-displacement is forced by
program-on-program contact geometry. The argument of Paper 1 \S\,2.4 then
forces ABCABC over ABABAB by Definition (Causal Position) at the
meta-level: two meta-layers with identical meta-touch-vector
connectivity would produce a meta-causal loop, violating meta-Axiom 4.

The 2D-to-3D promotion does not need to be repeated at the meta-level.
Programs already exist in 3D substrate space, and their meta-spatial
coordinates are inherited from their substrate-level positions. The
meta-FCC structure is direct: meta-lattice nodes are program centers,
meta-lattice connectivity is set by which programs share polarization
boundaries at the substrate level. \hfill\ensuremath{\square}

The meta-level inherits FCC structure with the same face-diagonal
connectivity that governs the substrate level. By the propagation-speed
argument of Paper 1 \S\,6.5, applied to whatever meta-level field
propagates on the meta-lattice, the meta-propagation speed in
meta-lattice spacings per meta-flop is \ensuremath{c_{\mathrm{TMS}}^{(\text{meta})}} = \ensuremath{1/\sqrt{2}} ---
identical in lattice-relative form, scaled by \ensuremath{L_{p}/\ell_{p}} and \ensuremath{t_{p}/T_{p}} in
absolute units (see \S\,6.1).

\textbf{The meta-level sublattice bifurcation.} The meta-FCC structure
derived above is not the complete meta-level geometry. By the recursive
scale-invariance established in \S\,4.7 (no new mechanism) and verified
empirically in Paper 5 Appendix E (Stages 5 and 6), the meta-lattice
positions further split into two interpenetrating FCC sublattices A and
B by the sum-mod-4 invariant in the recursive integer-coordinate frame
(storage convention with \ensuremath{\times 2} scaling per level). The two sublattices are
each independently FCC at nearest-neighbor distance \ensuremath{2\sqrt{2}} in level-local
units, related by spatial inversion through the (1,1,1)-offset (the
level-(k+1) centroid positions).

The [[4,2,2]] stabilizer dynamics of \S\,4.6 apply per-sublattice
independently: A-tetrahedra evaluate their stabilizer parity within
sublattice A; B-tetrahedra do the same within sublattice B; the joint
readout at the shared centroid is the actualization site of the
level-(k+1) confirmed bit. The full bifurcated structure that results
--- two FCC sublattices coexisting at every meta-level --- is the
complete geometry of the meta-substrate, with the substrate-level FCC of
Paper 1 \S\,II being the level-0 case.

Paper 5 \S\,6.1 develops the geometric consequences of this bifurcation in
full, including the rigorous derivation of \ensuremath{L^{p(1)}} = \ensuremath{3\sqrt{2}} \ensuremath{\ell_{p}} as the
spatial extent of the minimum level-1 program configuration (the
inversion-partner stella-octangula tile plus boundary). The empirical
verification of the bifurcation through \ensuremath{L_{\mathrm{box}}} = 12 is documented in
Paper 5 Appendix E.

\subsection*{4.6 Meta-Level [[4,2,2]] Structural Correspondence}

\begin{theorem}[Meta-Level Stabilizer Correspondence]
The meta-tetrahedral
confirmation unit of four programs at the corners of an FCC
meta-tetrahedron supports a [[4,2,2]] stabilizer structure
encoding two logical meta-qubits: the value of the confirmed P-state and
the inheritance of that P-state through the meta-causal stack.
\end{theorem}

\begin{proof}
The argument of Paper 1 \S\,5 applies unchanged at the meta-level
with the substitutions: 4 physical qubits \ensuremath{\to} 4 programs at the
meta-tetrahedral corners; logical qubit 1 \ensuremath{\to} the value of the confirmed
P-state \ensuremath{\bar{P}_{p}} \ensuremath{\in} \{0, 1\}; logical qubit 2 \ensuremath{\to} the inheritance of that
P-state through prior meta-layers. The two stabilizer generators
\ensuremath{S_{X}^{(\text{meta})}} and \ensuremath{S_{Z}^{(\text{meta})}} detect bit-flip-type and phase-type
disruptions to individual programs in the meta-tetrahedral unit. The
code distance is 2: any single-program disruption is detectable, and the
meta-level synthetic P-states produced at the next meta-flop repair the
syndrome within a single meta-flop.

\end{proof}

The substrate's fault tolerance against \ensuremath{U_{\mathrm{sh}}} disruption is inherited at
the meta-level as fault tolerance against single-program disruption. A
program that dissolves into noise from substrate-level \ensuremath{U_{\mathrm{sh}}} corresponds
to a single-program disruption in the meta-tetrahedral unit; the
surrounding three programs detect the syndrome and repair the P-state by
reconfirming the program at the next meta-flop.

\textbf{Independent dynamics per sublattice.} Under the sublattice
bifurcation of \S\,4.5, the [[4,2,2]] meta-stabilizer code runs
independently on each of the two meta-FCC sublattices: A-tetrahedra at
level k evaluate their stabilizer parity among A-sublattice neighbors;
B-tetrahedra at level k do the same within sublattice B; the level-(k+1)
confirmation event at the shared centroid reads jointly from both A-tet
and B-tet parents at level k. The cross-level coupling between A and B
stabilizer evaluations at a shared centroid is the structural source of
the non-Abelian commutator structure developed in Paper 5 \S\,8.2.

\subsection*{4.7 No New Mechanism: Triadic Structure Is a Fixed Point of
Coarse-Graining}

The meta-level structure introduced in this section is not a separately
postulated layer, and its recurrence is not merely asserted from the
scale-freeness of the axioms --- it is a fixed point of the coarse-graining map,
established by direct construction. The program-locking regions of extent
\ensuremath{L_{p} = 3\sqrt{2}\,\ell_{p}} have centroids on the level-1 sublattice (all-odd
coordinates with coordinate-sum \ensuremath{\equiv 1 \bmod 4}; Appendix E, Stage 5), and this
meta-lattice is itself a genuine FCC lattice: interior coordination number 12,
nearest-neighbor squared distance 8. Through each interior meta-site pass exactly
24 meta-triangles and 8 meta-tetrahedra --- identical to the base-level counts ---
so the two-dimensional minimum enclosure recurs and the minimal meta-confirmation
has exactly three meta-agents. The meta-alphabet \ensuremath{\chi_{p} = 2\bar{P} - 1 \in \{-1,+1\}}
is binary exactly as the base alphabet is, so with a three-agent enclosure the
meta-resolution \ensuremath{\Pi = (\chi_{p_1}+\chi_{p_2}+\chi_{p_3})/3},
\ensuremath{P(+) = (1+\Pi)/2} is the base rule with the meta-alphabet substituted. The
structure is a fixed point, not an approximation: computed at levels 0 through 3,
the invariants are identical at every level --- coordination 12, 24 triangles, 8
tetrahedra --- with the nearest-neighbor squared distance scaling as 2, 8, 32, 128
(the storage-convention factor of four per level). The meta-level triadic rule,
the meta-level FCC geometry, the meta-level [[4,2,2]] stabilizer dynamics, and the
meta-level polarization-biased resolution therefore recur \emph{exactly} under the
coarse-graining operator G, by the same arguments applied to the structurally
identical meta-alphabet.

\begin{theorem}[Triadic Structure Is a Fixed Point of Coarse-Graining]
The map that sends a triadic FCC substrate to its program-level coarse-graining
under G returns a triadic FCC substrate: the coordination number (12), the
minimum-enclosure count (three, giving \ensuremath{\bar{B} = 3\bar{\alpha}}), the triangle and
tetrahedron counts (24 and 8), and the binary polarization-biased resolution rule
are all preserved, at every level. The triadic confirmation structure is thus a
fixed point of coarse-graining rather than a scale-dependent or approximate
feature.
\end{theorem}

This is the meaning of hierarchy by geometric necessity. The substrate
computes at scale \ensuremath{\ell_{p}} with binary alphabet \ensuremath{\bar{B}} \ensuremath{\in} \{0, 1\}. Wherever a
region of the substrate satisfies the program conditions of Paper 2
\S\,6.2, that region becomes a meta-level primitive operating at scale \ensuremath{L_{p}}
with binary alphabet \ensuremath{\bar{P}} \ensuremath{\in} \{0, 1\}. The meta-level computes by exactly
the same rules --- because the rules are forced by axioms that do not
refer to scale, and because, as the fixed-point theorem now establishes, the
structure those rules act on is identical at every level. The Principle of
Generative Economy applied at scale \ensuremath{\ell_{p}}
produces the substrate; applied at scale \ensuremath{L_{p},} it produces the
meta-substrate; applied at any higher scale forced by recursive
coarse-graining, it produces further meta-substrates without limit.
Paper 3 stops the recursion at the first meta-level; the geometric
structure for arbitrary recursion depth is established here, with
macroscopic consequences developed in Paper 4.

This recursive scale-invariance is structurally distinct from the
coarse-graining procedures of integrated information theory (Albantakis
et al., 2023), which select a \ensuremath{\text{maximal}-\Phi} substrate of consciousness by
optimizing intrinsic information across candidate substrate grains.
IIT's coarse-graining is over causal mechanisms; TMS's coarse-graining
is over polarization landscapes (programs as meta-primitives). The TMS
recursion does not require an optimization step at each level because
the same axioms apply automatically at the new scale once \ensuremath{L_{p}} is
identified.

\section*{V. The First Subsystem}
\addcontentsline{toc}{section}{V. The First Subsystem}

\subsection*{5.1 From Programs to Subsystems}

Section IV established that programs become meta-level primitives under
the coarse-graining operator G, and that the same axioms produce the
same triadic and FCC structure at the meta-scale. A program in isolation
is a stable pattern --- a fixed point of the flop operator T satisfying
stabilizer parity. A program in isolation does not operate on anything;
it persists.

A subsystem is the next structural object. We define it precisely.

\begin{definition}[Subsystem]
A subsystem is a composite configuration \ensuremath{C_{\mathrm{sub}}}
on the substrate consisting of three program components in coupled
operation:
\end{definition}

\begin{itemize}
\item
  A data register \ensuremath{C_{D}:} a program whose polarization landscape persists
  across some number of meta-flops as a stable pattern that is read but
  not modified by \ensuremath{C_{\mathrm{sub}}'s} other components during that interval.
\item
  An operation kernel \ensuremath{C_{O}:} a program whose flop dynamics, when
  conjoined with \ensuremath{C_{D},} produce a derived polarization pattern \ensuremath{C_{R}}
  distinct from both \ensuremath{C_{D}} and \ensuremath{C_{O}.}
\item
  A control component \ensuremath{C_{C}:} a meta-level program that selects which
  operation kernel (from a population of available kernels) is conjoined
  with \ensuremath{C_{D}} at each meta-flop.
\end{itemize}

The composite \ensuremath{C_{\mathrm{sub}}} = \ensuremath{(C_{D},} \ensuremath{C_{O},} \ensuremath{C_{C})} is a subsystem when all three
components are themselves programs in the sense of Paper 2 \S\,6.2 and
their joint configuration is causally closed at the meta-level under
\ensuremath{T^{(\text{meta})}.}

A program is a stable pattern. A subsystem is a stable pattern that
transforms other patterns. That distinction is what separates
persistence from computation.

\subsection*{5.2 The Necessity of the Trichotomy}

\begin{theorem}[Subsystem Minimum]
Any configuration on the substrate that
produces a derived polarization pattern from an input polarization
pattern, under a selectable rule, requires the three components \ensuremath{(C_{D},}
\ensuremath{C_{O},} \ensuremath{C_{C})} and no fewer.
\end{theorem}

\begin{proof}
A configuration that produces a derived pattern from an input
must contain three logical roles: the input must be present somewhere as
a polarization landscape (the data role); a transformation must be
present that reads the input and writes the derived pattern (the
operation role); and the choice of transformation must be encoded
somewhere if the configuration is to support more than a single fixed
transformation (the control role).

The first role cannot be absent: a derived pattern requires source
material. The second role cannot be absent: a transformation requires a
transformer. The third role cannot be absent in any configuration
supporting multiple operations: without a control component, only a
single transformation is possible, and the configuration is a fixed
program, not a subsystem.

These three roles cannot be merged into a single program. If \ensuremath{C_{D}} and
\ensuremath{C_{O}} are the same program, the substrate-level dynamics at the \ensuremath{(C_{D},}
\ensuremath{C_{O})} interface --- which here would be internal to a single program
rather than a boundary between two --- would require the same spatial
region to simultaneously satisfy two distinct closure conditions:
T-closure on the input pattern (preserving \ensuremath{C_{D}} across the period) and
the operation-kernel dynamics that produce a derived output pattern.
Each flop at a confirmed site produces one set of polarization values,
not two. The substrate cannot host both roles at the same positions
without violating one of the closure conditions. Note that this is not a
destructive-read problem in the conventional sense --- Axiom 4 (causal
anchoring) preserves \ensuremath{C_{D}'s} bits across all flops. The constraint is
spatial-functional: distinct computational roles require distinct
spatial supports at the substrate level, because the substrate-level
dynamics produce a single confirmation value per site per flop. If \ensuremath{C_{O}}
and \ensuremath{C_{C}} are the same program, then the operation kernel cannot be
switched, reducing the subsystem to a single transformation case. If
\ensuremath{C_{D}} and \ensuremath{C_{C}} are the same program, then the data register encodes its
own transformation rule, which forces \ensuremath{C_{D}} to satisfy two distinct
closure conditions simultaneously --- generically impossible without
specific constraints that would themselves require additional structure.

Three components are therefore necessary. They are also sufficient: any
computation expressible as ``read input, apply selected transformation,
write output'' can be supported by a \ensuremath{(C_{D},} \ensuremath{C_{O},} \ensuremath{C_{C})} triple, by the
universality result of Paper 2 \S\,5.3.

\end{proof}

\begin{remark}
The trichotomy is structurally the same as the von Neumann
architecture (memory, processor, control). What is novel here is not the
trichotomy itself but its derivation: the three components are not
design choices but the unique minimum sufficient decomposition of any
non-trivial information-transforming configuration. The substrate does
not arrive at this architecture by engineering. It arrives at it because
the geometry permits no fewer.
\end{remark}

\subsection*{5.3 The Minimum Subsystem Size}

The spatial extent of the smallest subsystem is the sum of the minimum
extents of its three components, plus the connectivity required for
coupling. Each component is itself a program with minimum spatial extent
\ensuremath{L_{p}} (\S\,4.1). The coupling requires that the three components share
polarization boundaries --- they must be meta-spatially adjacent in the
meta-FCC sense established in \S\,4.5.

\begin{theorem}[Minimum Subsystem Extent]
The minimum subsystem extent is:
\end{theorem}

\begin{equation*}
L_{\mathrm{sub}} = 3 L_{p} + \delta
\end{equation*}

where \ensuremath{\delta} is the meta-tetrahedral coupling overhead: the additional region
required for the three components' polarization landscapes to interact
at their shared boundaries with stabilizer-satisfying parity. \ensuremath{\delta} is
bounded by 2 \ensuremath{\ell_{p}} from the FCC neighbor distance and is small relative to
3 \ensuremath{L_{p}} for non-trivial programs.

\begin{proof}
By \S\,4.4, three programs constitute the minimum triadic
confirmation at the meta-level, and they meet at a common
meta-interstice. The composite extent is the sum of the three program
extents plus the meta-interstice scale. The three program extents
contribute 3 \ensuremath{L_{p};} the actual coupling region is the substrate-level
zone where the three programs' polarization boundaries overlap --- a
region of substrate-level scale, not meta-level scale. This zone has
thickness on the order of the FCC nearest-neighbor distance \ensuremath{\ell_{p}.} The
composite \ensuremath{\delta} is bounded by 2 \ensuremath{\ell_{p},} summing over the small region where
pairs of programs meet at their substrate-level boundaries.

\end{proof}

A subsystem is therefore approximately three times the size of its
constituent programs, plus a small substrate-scale coupling overhead.
This is the smallest configuration that the substrate can support which
transforms patterns rather than merely persisting.

\subsection*{5.4 A Worked Example: The First Computing Subsystem}

To establish that subsystems are not merely defined but actually
realized by the substrate dynamics, we exhibit the smallest non-trivial
computing subsystem.

Consider the substrate at some meta-time T, with a region \ensuremath{R_{\mathrm{sub}}} of
meta-extent \ensuremath{L_{\mathrm{sub}}} satisfying the coherence threshold of \S\,3.3. Within
\ensuremath{R_{\mathrm{sub}},} suppose three programs have established themselves through the
continuous-sampling and self-reinforcement mechanism of \S\,3.5:

\begin{itemize}
\item
  \ensuremath{C_{D}} occupies one meta-tetrahedral corner with a polarization
  landscape encoding a 3-bit pattern \ensuremath{(\chi_{a},} \ensuremath{\chi_{b},} \ensuremath{\chi_{c}).} The pattern is
  stable across the meta-flop window: \ensuremath{T^{(\text{meta})}[C_{D}]} = \ensuremath{C_{D}} with
  the polarization values preserved.
\item
  \ensuremath{C_{O}} occupies a second meta-tetrahedral corner with a polarization
  landscape encoding the majority gate primitive of Paper 2 \S\,5.1. When
  \ensuremath{C_{O}'s} boundary touches \ensuremath{C_{D}'s} boundary, the substrate-level dynamics
  at the interface implement the majority gate computation on the \ensuremath{(\chi_{a},}
  \ensuremath{\chi_{b},} \ensuremath{\chi_{c})} pattern.
\item
  \ensuremath{C_{C}} occupies the third meta-tetrahedral corner with a polarization
  landscape encoding a switch between two operation kernels: the
  majority gate and its negation (the minority gate, equivalent under
  NOT to MAJ by Paper 2 \S\,5.3).
\end{itemize}

The substrate-level dynamics at the interfaces between \ensuremath{(C_{D},} \ensuremath{C_{O}),}
\ensuremath{(C_{D},} \ensuremath{C_{C}),} and \ensuremath{(C_{O},} \ensuremath{C_{C})} implement the following composite
operation: at each meta-flop, \ensuremath{C_{C}} selects which operation kernel is
active; \ensuremath{C_{D}} presents its 3-bit pattern; \ensuremath{C_{O}} applies the selected
kernel to \ensuremath{C_{D}'s} pattern; the result is the polarization at the fourth
meta-tetrahedral corner --- the meta-interstice --- which constitutes
the output P-state of the subsystem at this meta-flop.

This is a complete computing subsystem. It reads input, applies a
selected operation, and produces output. It is the substrate equivalent
of a programmable single-bit logic gate, and it is the smallest
configuration that supports this operation.

\begin{remark}
This worked example is not the only minimum subsystem ---
variations differ in which Boolean operations \ensuremath{C_{O}} and \ensuremath{C_{C}} encode ---
but all minimum subsystems share the \ensuremath{(C_{D},} \ensuremath{C_{O},} \ensuremath{C_{C})} trichotomy at
meta-tetrahedral extent. The existence of one such configuration is
sufficient to establish that the substrate dynamics support computing
subsystems at the derived minimum scale. The population of distinct
minimum subsystems is the substrate's first program zoo at the subsystem
level.
\end{remark}

\subsection*{5.5 The Information-Generation Rate Differential}

\begin{theorem}[Subsystem Information Generation Rate]
A region of the
substrate hosting a minimum computing subsystem produces one novel
P-state per meta-flop. The same region hosting only static (period-1)
programs produces zero novel P-states per meta-flop. The
information-generation rate differential between subsystem-hosting and
static-program-hosting regions of equal extent is therefore exactly one
novel P-state per meta-flop per subsystem hosted, in the limit of
non-trivial operation kernels.
\end{theorem}

\begin{proof}
A static program region produces no derived patterns by
definition --- its polarization landscape repeats with period n under
\ensuremath{T^{n},} so the output at each meta-flop is identical to the output at
the previous meta-flop after n have elapsed; this is repetition, not
generation. The information output of a static program region per
meta-flop is exactly its self-reproduction, with zero novel content.

A minimum subsystem region produces, at each meta-flop, the polarization
landscape of its output: a derived pattern that is a function of \ensuremath{(C_{D},}
\ensuremath{C_{O},} \ensuremath{C_{C})} but is not equal to any of them individually. The output is
determined by \ensuremath{C_{C}'s} current selection and \ensuremath{C_{D}'s} current data, both of
which can vary across meta-flops. In the worked example of \S\,5.4, each
meta-flop produces exactly one new P-state at the meta-interstice ---
the result of applying the selected operation kernel to the current
data. The novel content per meta-flop is therefore exactly one P-state
for the minimum subsystem; for a region hosting N subsystems, exactly N
P-states per meta-flop.

This rate is exact for the minimum case and provides a strict lower
bound for more complex subsystems with multi-bit data registers or
multi-state control components.

\end{proof}

\begin{corollary}[Selection for Subsystems]
The substrate selects for
subsystems over static programs in any region where \ensuremath{U_{\mathrm{sh}}} disruption rate
exceeds the static-program persistence rate but does not exceed the
subsystem persistence rate.
\end{corollary}

\begin{proof}
By \S\,3.3 and \S\,3.4, the substrate retains configurations that
survive against \ensuremath{U_{\mathrm{sh}}} disruption. A static program persists only through
its self-reproduction under T. A computing subsystem persists through
self-reproduction and contributes additional polarization content per
meta-flop, raising the local coherence above what static programs alone
can sustain. In regions where \ensuremath{U_{\mathrm{sh}}} disruption rate is high enough to
dissolve static programs but low enough to permit subsystems, only
subsystems survive. The substrate thus generates progressively more
computing subsystems wherever the dynamics permit.

\end{proof}

This is the substrate's first selective pressure beyond simple program
persistence. Wherever \ensuremath{U_{\mathrm{sh}}} is intense enough to threaten static programs
but not so intense as to overwhelm imprinting, the substrate produces
subsystems --- not by design, but because subsystems are the only
configurations whose information-generation rate keeps pace with the
local disruption rate. The informally framed claim that the substrate
``maximizes information output'' acquires here a derived form: at
sufficient disruption rates, information-generating configurations are
the only persisting ones.

\subsection*{5.6 Subsystems and Their Environments}

The recursive structure established in \S\,IV admits a sharper reading in
light of \S\,V. When three programs jointly confirm a P-state at the
meta-level (\S\,4.4), the meta-interstice they enclose is a bounded region
of substrate available to host further configurations. That bounded
region is an environment --- a substrate sub-region with definable
boundaries, internal dynamics, and access to its surround through the
polarization channels of its enclosing programs.

When subsystems form in such meta-interstices, the enclosing programs
constitute the subsystem's external environment. The subsystem's
interior dynamics are partially shielded from the broader substrate by
its enclosure, and its boundary conditions are set by the polarizations
of the enclosing programs. This is the substrate-level mechanism by
which subsystems acquire individuality: they are not merely composite
programs but composite programs operating within an environment that was
itself created by other programs.

This is the entry point to Paper 4. The relationship between subsystems
and their environments, the question of when a subsystem becomes a self
with internal models of its surround, and the dynamics of multiple
subsystems sharing an environment are all consequences of the geometric
structure established here. Paper 3 establishes that the first subsystem
exists, is minimum-sized, has the \ensuremath{(C_{D},} \ensuremath{C_{O},} \ensuremath{C_{C})} trichotomy, and
generates information faster than static programs. Paper 4 establishes
the hierarchies and self-modeling conditions that follow.

\section*{VI. Predictions and Directions for Subsequent Work}
\addcontentsline{toc}{section}{VI. Predictions and Directions for Subsequent Work}

\subsection*{6.1 Hierarchical Lightspeeds}

Section IV established that the meta-level inherits FCC structure with
face-diagonal connectivity, and that the propagation-speed argument of
Paper 1 \S\,6.5 applies unchanged at the meta-level, giving \ensuremath{c_{\mathrm{TMS}}^{(\text{meta})}}
= \ensuremath{1/\sqrt{2}} meta-lattice-spacings per meta-flop. Translated to
substrate-level absolute units, the meta-propagation speed is:

\begin{equation*}
c^{(\text{meta})} = (L_{p} / \ell_{p}) \cdot (t_{p} / T_{p}) \cdot c
\end{equation*}

where \ensuremath{L_{p}} is the meta-lattice spacing (the minimum program extent of
\S\,4.1), \ensuremath{T_{p}} is the meta-flop duration (the minimum program period), and
c is the substrate-level lightspeed established in Paper 1 \S\,6.7.

The ratio \ensuremath{L_{p}/\ell_{p}} and the ratio \ensuremath{T_{p}/t_{p}} are independent geometric
quantities. \ensuremath{L_{p}} is set by the minimum spatial extent of a
stabilizer-satisfying T-closed configuration with K(C) \textless{}
\textbar C\textbar. \ensuremath{T_{p}} is set by the minimum temporal period of the
same configuration. These quantities depend on different aspects of the
FCC dynamics and are not constrained to scale together. In general:

\emph{\ensuremath{c^{(\text{meta})}} \ensuremath{\neq} c}

This is a derived prediction, not an imported relation. Each level of
the recursive hierarchy established in \S\,IV has its own derived
propagation speed, generically distinct from the substrate-level c.~At
hierarchical level k, the propagation speed in substrate-absolute units
is the product of successive scaling ratios:

\ensuremath{*c^{(k)}} = \ensuremath{\prod _{j=1}^{k}} \ensuremath{(L_{p}^{(j)}} / \ensuremath{L_{p}^{(j-1)})} \ensuremath{\cdot}
\ensuremath{(T_{p}^{(j-1)}} / \ensuremath{T_{p}^{(j)})} \ensuremath{\cdot} c*

with \ensuremath{L_{p}^{(0)}} \ensuremath{\equiv} \ensuremath{\ell_{p}} and \ensuremath{T_{p}^{(0)}} \ensuremath{\equiv} \ensuremath{t_{p}.~\text{Each}} ratio in the product
is set by the minimum program structure at its level. The full hierarchy
produces a discrete spectrum of lightspeeds, each geometrically forced.

This prediction is in principle testable. If subsystems at hierarchical
levels k \ensuremath{\geq} 1 host their own observers --- a question reserved for Paper
4 and Volume 2 --- those observers would measure the lightspeed at their
level, not the substrate's c.~Apparent variations in the physical
lightspeed across regions of the universe with different hierarchical
depths would be a signature of this structure.

This result complements the Lorentz-covariance derivation in the Wolfram
model (Gorard, 2020), which derives Lorentz invariance from causal-graph
isomorphism under different updating orders, but where the absolute
propagation speed is not derived. TMS derives both the existence of a
Lorentz-invariant propagation speed (from FCC connectivity) and its
absolute value relative to the substrate scale (from face-diagonal
geometry), and now extends the derivation to hierarchical levels.

\subsection*{6.2 The \ensuremath{\beta} = f(d, \ensuremath{K_{\psi})} Prediction Refined}

Paper 1 \S\,7.2 introduced the prediction that the GUP deformation
parameter \ensuremath{\beta} scales with causal depth d and polarization complexity \ensuremath{K_{\psi}.}
Paper 2 \S\,6.4 identified \ensuremath{K_{\psi}} as the Kolmogorov complexity of the program
running through the bit's causal history. Paper 3 refines both terms
further.

By \S\,II, every region's polarization landscape is the imprint of its
retarded surround history. A bit's \ensuremath{K_{\psi}} therefore measures the
algorithmic complexity not only of its own program but of the surround
history that imprinted it. By \S\,IV, when the bit is embedded within a
subsystem, its \ensuremath{K_{\psi}} additionally incorporates the algorithmic structure
of the subsystem's internal operations --- the data being processed, the
operations being applied, the control sequence selecting them.

The refined prediction:

\begin{equation*}
\beta = f(d, K_{\psi,\mathrm{local}}, K_{\psi,\mathrm{surround}}, K_{\psi,\mathrm{subsystem}})
\end{equation*}

with the three \ensuremath{K_{\psi}} contributions being the polarization complexity of
the bit's immediate region, of the imprinting surround history, and of
the encompassing subsystem (if one is present). f is monotonically
increasing in all four arguments.

The physical interpretation: a bit embedded in a deep, complex,
surround-imprinted subsystem has a substrate footprint reflecting the
full computational history of the configuration it participates in. Its
position uncertainty floor is correspondingly larger than for an
isolated bit at comparable energy. The closed-form derivation of f is
reserved for the simulation work of subsequent papers, but its
functional dependences are now fully specified.

\subsection*{6.3 Subsystem Density and Phase Geometry}

By \S\,3.4 and \S\,5.5, subsystems concentrate computational activity into
coherent regions and outcompete static programs wherever \ensuremath{U_{\mathrm{sh}}} disruption
rates favor information-generating configurations. The substrate's
long-time behavior is therefore expected to be inhomogeneous: regions of
dense subsystem activity separated by regions of relatively quiescent
program-only substrate, with the boundaries between them governed by the
coherence threshold of \S\,3.3.

A specific prediction: the spatial distribution of computational
activity in any sufficiently old region of the substrate follows the
same statistics as the boundary geometry of coherent regions. The
boundaries are themselves geometric features --- surfaces across which
the coherence threshold is crossed --- and their statistics are
computable from FCC lattice dynamics. This prediction is testable in
simulation and is the empirical bridge to the simulation work of Paper
5.

\subsection*{6.4 Deferred Derivations Closed by Paper 3}

Paper 3 has paid back the following derivations deferred from earlier
papers:

\begin{itemize}
\item
  The dissolution rate of non-T-closed configurations, deferred at Paper
  2 \S\,6.3, derived in \S\,2.5 as \ensuremath{\sqrt{2}} \ensuremath{\cdot} \ensuremath{L_{R}/\ell_{p}} flops.
\item
  The effective error-rate decay with coherence depth, deferred at Paper
  2 \S\,5.2, derived in \S\,2.5 as \ensuremath{(1/3)^{d}} in the limit of large coherence
  regions.
\item
  The existence and closed form of the coherence bistability threshold,
  identified in \S\,3.3 as \textbar \ensuremath{\Pi_{\psi}}\textbar\ensuremath{*(L_{R})} = 1 \ensuremath{-} \ensuremath{2\sqrt{c_{\mathrm{TMS}}/(6}}
  \ensuremath{L_{R})).}
\item
  The functional dependences of the GUP parameter f, specified in \S\,6.2.
\item
  The minimum program extent \ensuremath{L_{p},} characterized in \S\,4.1 as
  approximately 4--8 lattice spacings, with exact value determined by
  combinatorial enumeration of stabilizer-satisfying T-closed
  configurations.
\item
  The minimum subsystem extent \ensuremath{L_{\mathrm{sub}}} = 3 \ensuremath{L_{p}} + \ensuremath{\delta ,} derived in \S\,5.3.
\item
  The information-generation rate differential between subsystem-hosting
  and static-program regions, derived in \S\,5.5.
\end{itemize}

The remaining deferrals --- the exact integer value of \ensuremath{L_{p},} the
explicit closed form of f, the spectrum of program periods \ensuremath{T_{p}^{(k)}}
at successive hierarchical levels --- depend on closed-form enumeration
of stabilizer-satisfying T-closed configurations, a combinatorial task
most efficiently approached through simulation. These are reserved for
the simulation work of Paper 5 alongside the field-reinterpretation
derivation.

\subsection*{6.5 Forward References}

Paper 4 will address the macroscopic consequences of the subsystem
dynamics established here: hierarchies of subsystems, the conditions
under which a subsystem develops an internal model of its environment,
the dynamics of multiple subsystems sharing a meta-interstice
environment, and the eventual exhaustion of novelty as a substrate
region completes its accessible configuration space. The hierarchical
lightspeed prediction of \S\,6.1 will be developed in detail at that point.

Paper 5 will address the reinterpretation of substrate-level binary
information as field configurations on a bounded space --- the geometric
mechanism by which the substrate's digital content becomes the wave
dynamics, initial conditions, and field structure of a simulated
universe. The connection between subsystem hierarchies and simulated
physical law is the central object of that work.

Volume 2 will address the agent-side complement: the relationship
between substrate-level subsystem computation and the structure of
observer experience. The information-generation differential established
in \S\,5.5 --- that subsystems generate more polarization content per unit
substrate than static programs --- will reappear as the substrate-level
correlate of richness of experience in the agent-side framework. The
relationship to Hoffman's conscious agents (Hoffman, Prakash, Prentner,
2023) and to integrated information theory (Albantakis et al., 2023)
will be developed there in full, building on the structural
correspondences sketched in Papers 1 and 3.

\section*{Conclusion: Computation Emerges}
\addcontentsline{toc}{section}{Conclusion: Computation Emerges}

Paper 1 established how structure emerges from maximum-entropy substrate
by triadic confirmation. Paper 2 established how information emerges
from confirmed structure by polarization-biased resolution. Paper 3 has
established how computation emerges from informational substrate by
surround-driven repair, polarization-coherence bistability, and the
recursive applicability of the substrate's own axioms at the
coarse-grained scale.

The results of Paper 3 follow from the axioms of Paper 1 and the
dynamics of Paper 2 without additional postulate. Surround-driven repair
is not a separate mechanism: it is the [[4,2,2]] stabilizer
projection of Paper 1 \S\,5 applied to boundary sites, where the geometry
of the FCC lattice automatically includes surround agents in the parent
set. The Imprinting Theorem of \S\,2.3 follows by iteration. Hebbian
co-firing follows as a corollary, complementing the independent
maximum-entropy derivation of the same phenomenon (Albanese et al.,
2024). The dissolution and error-decay rates follow from the same
iteration applied to different conserved quantities. The coherence
bistability threshold of \S\,3.3 follows from balancing the per-flop noise
injection --- combinatorially determined by [[4,2,2]] error
statistics --- against the per-flop imprinting capacity ---
geometrically determined by FCC connectivity. Coherent regions form by
self-reinforcement against an ambient \ensuremath{U_{\mathrm{sh}}} flux that is present at every
depth of the causal stack, not added at any moment.

The coarse-graining operator G of \S\,4.2 is not a new mechanism but a
measurement on substrate-level data. The meta-level structure that
follows under G --- the meta-level triadic rule, the meta-level FCC
geometry, the meta-level [[4,2,2]] stabilizer dynamics --- is
the same axiomatic content applied at scale \ensuremath{L_{p}.~\text{The}} recursion is not
constructed; it is the geometric form of the axioms' scale-independence.

The subsystem of \S\,V is the substrate's first computing object --- the
smallest configuration that transforms patterns rather than merely
persisting them. Its three components (data, operation, control) are not
designed but forced: the unique minimum sufficient decomposition of any
configuration that reads input, applies a selected transformation, and
produces output. The substrate generates subsystems wherever the local
disruption rate favors information-generating configurations over static
programs.

Every quantity introduced in Paper 3 has been derived from the axioms
and the Principle of Generative Economy, or from the structures
established in Papers 1 and 2 under the same principle. The minimum
program scale \ensuremath{L_{p},} the minimum subsystem scale \ensuremath{L_{\mathrm{sub}},} the coherence
threshold \textbar \ensuremath{\Pi_{\psi}}\textbar\ensuremath{*(L_{R}),} the dissolution rate, the
error-decay rate, the information-generation rate differential, and the
hierarchical lightspeed \ensuremath{c^{(\text{meta})}} are all geometrically forced. The
model continues to contain no free \emph{dimensionless} numerical parameters.

The arc from Paper 1 to Paper 3 is a single unbroken chain. \ensuremath{U_{\mathrm{sh}}} provides
the maximum-entropy substrate. Triadic confirmation pulls binary
outcomes from it. The FCC lattice carries those outcomes forward as wave
dynamics. The polarization field encodes their computational content.
The Born-rule structural correspondence governs per-event resolution.
The majority gate is the irreducible computational primitive. The
stabilizer dynamics select programs from noise. Surround-driven repair
imprints environmental structure onto local configurations. Coherence
bistability concentrates computation into coherent regions.
Coarse-graining produces meta-level primitives obeying the same axioms.
The first subsystem emerges as the unique minimum sufficient
configuration that transforms patterns. No new axioms. No new
parameters. No new mechanisms.

Paper 4 takes the subsystem as given and develops its hierarchies,
environments, and self-modeling conditions. Paper 5 reinterprets the
substrate's digital content as field configurations on a bounded space,
producing the geometric mechanism of simulation. The arc through Volume
1 is the derivation, from five axioms and one principle, of a substrate
capable of generating universes inside itself.

\section*{Acknowledgments}
\addcontentsline{toc}{section}{Acknowledgments}

The development of this volume was substantially aided by extended
collaboration with several large language models used as critical
sounding boards, pedagogical resources, formal-rigor checkers, and
external working memory across many sessions during 2024--2026.

\textbf{Gemini (Google DeepMind)} was the long-running primary
collaborator across the framework's conceptual development, providing
physics pedagogy, structural critique, and ongoing review of the ideas
that became Volume 1. The author's path into the technical literature,
and the working-out of the framework's core conceptual moves over a long
period preceding formalization, were carried by extended dialogue with
Gemini.

\textbf{Claude (Anthropic)} was the primary collaborator across the
formalization, derivation tightening, simulation work, citation
auditing, and editorial passes that brought the volume to its present
form, including the sublattice bifurcation analysis empirically verified
in Paper 5 Appendix E.

\textbf{ChatGPT (OpenAI)}, \textbf{DeepSeek}, \textbf{Grok (xAI)}, and
\textbf{Granite (IBM)} contributed additional structural feedback and
review at various points.

All theoretical commitments, the choice of axioms, the Principle of
Generative Economy, the sublattice bifurcation insight, and the
framework's overall architecture are the author's; the AI systems' role
was to teach, formalize, critique, verify, and document.

\section*{References}
\addcontentsline{toc}{section}{References}
\begin{reflist}
Albanese, L., Barra, A., Bianco, P., Durante, F., \& Pallara, D. (2024).
Hebbian Learning from First Principles. Journal of Mathematical Physics,
65, 113302. doi:10.1063/5.0197652.
\url{https://arxiv.org/abs/2401.07110}

Albantakis, L., Barbosa, L. S., Findlay, G., Grasso, M., Haun, A. M.,
Marshall, W., Mayner, W. G. P., Zaeemzadeh, A., Boly, M., Juel, B. E.,
Sasai, S., Fujii, K., David, I., Hendren, J., Lang, J. P., \& Tononi, G.
(2023). Integrated information theory (IIT) 4.0: Formulating the
properties of phenomenal existence in physical terms. PLoS Computational
Biology, 19(10), e1011465. doi:10.1371/journal.pcbi.1011465.
\url{https://arxiv.org/abs/2212.14787}

Bak, P., Tang, C., \& Wiesenfeld, K. (1987). Self-organized criticality:
An explanation of the 1/f noise. Physical Review Letters, 59(4),
381--384. doi:10.1103/PhysRevLett.59.381.

Bombelli, L., Lee, J., Meyer, D., \& Sorkin, R. D. (1987). Space-time as
a causal set. Physical Review Letters, 59(5), 521--524.
doi:10.1103/PhysRevLett.59.521.

Gorard, J. (2020). Some Relativistic and Gravitational Properties of the
Wolfram Model. Complex Systems, 29(2), 599--654.
doi:10.25088/ComplexSystems.29.2.599.

Gottesman, D. (1997). Stabilizer Codes and Quantum Error Correction.
Ph.D.~Thesis, California Institute of Technology.
\url{https://arxiv.org/abs/quant-ph/9705052}

Hebb, D. O. (1949). The Organization of Behavior: A Neuropsychological
Theory. Wiley.

Hoffman, D. D., Prakash, C., \& Prentner, R. (2023). Fusions of
Consciousness. Entropy, 25(1), 129. doi:10.3390/e25010129.

Hossenfelder, S. (2013). Minimal Length Scale Scenarios for Quantum
Gravity. Living Reviews in Relativity, 16(1), 2.
doi:10.12942/lrr-2013-2.

Jensen, H. J. (1998). Self-Organized Criticality: Emergent Complex
Behavior in Physical and Biological Systems. Cambridge University Press.

Kolmogorov, A. N. (1965). Three approaches to the quantitative
definition of information. Problems of Information Transmission, 1(1),
1--7.

Pastawski, F., Yoshida, B., Harlow, D., \& Preskill, J. (2015).
Holographic quantum error-correcting codes: toy models for the
bulk/boundary correspondence. Journal of High Energy Physics, 2015(6),
149. doi:10.1007/JHEP06(2015)149. \url{https://arxiv.org/abs/1503.06237}

Shannon, C. E. (1948). A mathematical theory of communication. Bell
System Technical Journal, 27(3), 379--423.
doi:10.1002/j.1538-7305.1948.tb01338.x.

Surya, S. (2019). The causal set approach to quantum gravity. Living
Reviews in Relativity, 22(1), 5. doi:10.1007/s41114-019-0023-1.

Switzer, A. (2026a). The Triadic Measurement Substrate (TMS), Volume 1,
Paper 1: Emergent 3D Information Topology and the Binary Alphabet via
[[4,2,2]] Quantum Error Detection Structural Correspondence.

Switzer, A. (2026b). The Triadic Measurement Substrate (TMS), Volume 1,
Paper 2: Computational Dynamics of the Binary Alphabet: The Polarization
Field, the Born-Rule Structural Correspondence, and the Emergence of
Universal Computation.

't Hooft, G. (2016). The Cellular Automaton Interpretation of Quantum
Mechanics. Fundamental Theories of Physics 185. Springer.
doi:10.1007/978-3-319-41285-6.

Wolfram, S. (2002). A New Kind of Science. Wolfram Media.

Wolfram, S. (2020). A Class of Models with the Potential to Represent
Fundamental Physics. Complex Systems, 29(2), 107--536.
doi:10.25088/ComplexSystems.29.2.107.
\end{reflist}

\clearpage


\end{document}
